\subsection{Separating and testing residue planes}
\label{sec:plane-lifts}
\hypertarget{source-plane-lifts}{}

In the smaller-prime range of Section~\ref{sec:host-edge-tests}, the incoming lattices are indexed by residue planes. To add their ranks, each test must vanish to higher order on competing planes. Put \(r=r_G\).
An edge of type \(r\) has \(L_+/L_-\) of dimension \(r\) over
\(\mathbb F_q,\ q=p^{n'}\). For an edge of type \(r\), the exponent in the normalized pairing is
\[
                   A_r=\sum_s\sum_{i=1}^r\mu_{s,i}.
\] For a fixed source, competitors are
\(r\)-planes in \(p^{-1}L_-/L_-\). Choose coordinates taking one to the plane \(I_0=\{1,\ldots,r\}\) by an integral frame change.
At this plane the substitution multiplies the even vector
coordinates outside \(I_0\) and the odd covector coordinates
inside \(I_0\) by \(w\). Call these coordinates normal and the
others tangential. At a different plane use the corresponding
twisted change of basis before measuring divisibility.

\begin{proposition}[Neighbor tests on the smaller-prime range]\label{prop:residue-plane-neighbor-test}
Use the selected good cycle and the smaller-prime bounds of
Section~\ref{sec:host-edge-tests}, with its prescribed order
\(t_{\rm II}\), depth \(D_{\rm II}\), and cutoff
\(Y_2=x^{26}\). At each target let \(V_{\rm rec}\subseteq V\) be the subspace recovered from the incoming edges of type \(r\). Then
\[
 \dim V_{\rm rec}\ge(1-\kappa)\dim V,\qquad\kappa=O(Nx^{-12}).
\]

Here both endpoint orders are at most
\(\lceil t_{\rm II}\rceil\), and the evaluation on the recovered
directions has an integral right inverse after multiplication
by \(p^{D_{\rm II}}\).
\end{proposition}

The proof first works in a fixed plane. It separates that plane
from its competitors, lifts a space of residue symbols, and
constructs detecting covectors with the required vanishing.
Packing the resulting per-plane spaces at the same target
then gives the rank assertion.

\begin{lemma}[Frames, separation and normal directions on a good cycle]
The selected cycle admits \(r\)-row frames with all the required
Frobenius powers at most \(Q=p^{G'}\). In the adapted coordinates the separating polynomial vanishes on the other residue planes and is nonzero on the chosen plane, as specified in the proof. The bound for the Frobenius exponents \(p^{a_i}\) required in these frames is
\[
 p^{a_i}\le Q=p^{G'}\quad\text{for every required }i.
\]
 At a selected position the normal
coordinates on the chosen plane have integral polynomial lifts
after multiplication by the frame determinant and by \(w\),
to first order modulo \(w^2\).
\end{lemma}
\begin{proof}
In each marked block of Section~\ref{subsec:frobenius-selection}
choose the middle \(r\) of its \(3r\) isolated selected even
positions. These are consecutive tops. On either side there are
another \(r\) even tops in that block. Between the last position
of one frame and the first of the next there are consequently at
least \(2r\) tops. The first and last \(r\) have preceding gaps
greater than \(L_1\), by their isolation and the empty cells.
The first position of a frame has such a gap as well. From the
beginning of one frame to the end of the next the distance,
including both enclosing blocks, lies within \(G'\), even across
the partial block at wrap-around. Indeed marked block starts
are at distance at most \((M+2)l\), leaving two full block
lengths of slack in
\[
 G'=(M+5)l.
\]
The number of tops in
this segment is thus at most \(N_1\).

For any top after the
incoming frame, the \(r\)-th preceding even top is present in
the segment. Its distance is greater than \(L_1\), since
\(r=K+1\) and no cyclic interval of length \(L_2\ge2L_1\)
can contain \(r\) tops. In particular all powers taking any of
these \(r\) rows up to the tested top are at most \(Q=p^{G'}\).
These properties give the frame conditions and the twisting bounds used below.

\proofstep{Separate Pl\"ucker coordinates along the cycle.}
Write the frame positions as
\[
 a_{j,1}<\cdots< a_{j,r}=z_j
\]
in cyclic order, lifted to integers so that
\(z_{j+J}=z_j+n'\). For each \(j\), form the \(r\)-by-\(k\)
matrix of vector rows at \(a_{j,l}\), raising the \(l\)-th row's
entries to \(p^{z_j-a_{j,l}}\). Write \(P_{j,I}\) for the
minor on a subset \(I\) of size \(r\), \(D_j=P_{j,I_0}\) and
\[
 Q_j=p^{z_{j+1}-z_j}.
\]
There are \(r(k-r)\) subsets \(I\)
differing from \(I_0\) by one exchange. Use all of them in
\[
 R=D_0\prod_I\left(\prod_{j=0}^{J-1}D_j^{Q_j-1}
             -\prod_{j=0}^{J-1}P_{j,I}^{Q_j-1}\right),\quad
 R_0=\prod_j D_j^{\alpha_j},\quad
 \alpha_j=\mathbf1_{j=0}+r(k-r)(Q_j-1).
\]
Modulo \(w\) at a competing plane, let \(B\) be its rank-\(r\)
matrix of rows over \(\mathbb F_q\) in the distinguished frame
and \(b_I=\det B_I\). The \(j\)-th row matrix is a product of
an \(r\)-square matrix in the free variables with \(B^{[p^{z_j}]}\).

Indeed at a position \(a\) the surviving row is a linear
combination of the plane rows twisted by Frobenius at \(a\).
Powering moves them all to \(z_j\). Write Frobenius exponents
modulo \(n'\) in this residue calculation. Every minor thus has
the same polynomial factor at \(j\), times \(b_I^{p^{z_j}}\).
For \(b_I\ne0\) the factors telescope:
\[
 \prod_j (b_I^{p^{z_j}})^{Q_j-1}
    =b_I^{p^{z_J}-p^{z_0}}=1,\qquad z_J=z_0+n' .
\]
The final equality uses \(b_I\in\mathbb F_{p^{n'}}^\times\). If \(b_{I_0}=0\),
the initial \(D_0\) vanishes. Otherwise row reduce \(B\) in
the first \(r\) columns. In the resulting matrix \([I_r\mid C]\), a plane different
from the coordinate plane has some \(C_{ih}\ne0\). Exchanging column
\(i\) of \(I_r\) with column \(h\) of \(C\) gives
\[
 \det(e_1,\ldots,e_{i-1},C_{\bullet h},e_{i+1},\ldots,e_r)=C_{ih}\ne0 .
\]
Reordering the columns can only change the sign. The corresponding
difference factor therefore vanishes. This proves \(R\in wO[\mathbf X]\) after
substitution at every competitor.

At the chosen plane \(D_j\) depends only on tangential
coordinates. In every monomial of a one-exchange minor one
entry comes from a normal column. After substitution its
coefficient has \(w\)-order equal to its total normal degree.
Expanding \(R-R_0\), or \(R^c-R_0^c\), each time a one-exchange
term is chosen it costs at least \(Q_j-1\) normal degrees at
each \(j\). Thus, before collecting terms, each error term of normal degree
\(e\) has an explicit factor \(w^e\) and retains the tangential
factor
\[
 D_j^{\max(0,c\alpha_j-e)}
\]
at each frame for \(R^c\).
Here the word monomial refers to expansion after the undamaged
determinant factors have been set aside.

\proofstep{Recover a normal coordinate at the chosen plane.}
Use the \(r\) preceding selected even positions at each selected
\(s\), and power their vector rows to the index \(s\). Denote the
first \(r\) columns by \(U_s\) and put
\[
 T_s=\det U_s.
\]

Every power comes from a strictly earlier position. For even
\(s\) and \(h>r\), append column \(h\) and the current row to
form an \((r+1)\)-minor. At the distinguished plane its leading
expression is \(wT_s X_{s,h}\), with error divisible by \(w^2\).
Any other term using column \(h\) from an earlier row has at
least \(p\) powers of \(w\). At every plane the enlarged
matrix has rank at most \(r\) modulo \(w\), by the same
Frobenius factorization.

For odd \(s\), contract each powered
row with the current covector. At every plane these contractions
vanish modulo \(w\), since the surviving covector annihilates the
surviving row space. The leading terms at the distinguished plane
are the entries of
\[
 wU_s(X_{s,i})_{i\le r},
\]
again with errors
in \(w^2\). The adjugate of \(U_s\) expresses \(T_s X_{s,i}\)
as a tangential linear combination of their leading terms
divided by \(w\). This gives an order-one polynomial for each
required term. All letters are at selected positions.

\end{proof}

\paragraph{Correction to the prescribed precision.}

Separating one residue plane gives a congruence at the other planes. A correction is needed to reach the prescribed precision there while preserving the initial symbol.

\begin{lemma}[Correction on all residue planes]\label{lem:plane-normal-correction}
In this paragraph a factor
of order \(d\le2Q\) means an integral polynomial having order
at least \(d\) in every plane frame, with distinguished leading
expression \(S\) divisible by \(w^d\), homogeneous in the
normal coordinates of degree \(d\), and accurate modulo \(w^{d+1}\).
Use finitely many such factors \(K_l\) with expressions \(S_l\),
a power \(w^i\), \(i\ge0\), and any integral tangential polynomial
\(H\). Use only letters at tops. The level is \(L=i+\sum_l d_l\).

At any \(b\) with \(0\le b\le c\), consider the desired expression
\[
  w^i\prod_l S_l\, R_0^c\prod_jD_j^{b}
                  \prod_{s\ {\rm top}}T_s^{(2Q+3)b}\,H.
\]
In this expression all determinants are in the distinguished
coordinates. There is an integral polynomial whose substitution
agrees with it modulo \(w^{L+b}\) and whose order at competitors
is at least \(L+b\). For an initial level \(L\ge o\), taking
\(b=c\) gives the precision \(o+c\) needed before division by
\(w^o\). The construction can preserve any prescribed full
multidegree of the desired expression.
\end{lemma}

\begin{proof}
To prove the lemma, form the raw product by using the true
factors \(K_l,R\) and the true determinants and tangential
polynomial in the distinguished frame. The latter determinants
and polynomial restrict exactly there, since their entries are
tangential. Every term of the raw product has the required
competitor divisibility as a polynomial, by the universal
orders and the factor \(R^c\). The case \(b=0\) also follows
by taking zero. For \(b>0\), expand its error at the designated
plane. A replaced factor \(K_l\) contributes a monomial of
normal degree \(m\) with coefficient divisible by
\(w^{d_l+e_l}\), where
\[
 e_l=\max(1,m-d_l).
\]

Indeed the substitution of an integral polynomial first incurs
\(w^m\) on that monomial, and its difference from \(S_l\)
is divisible by \(w^{d_l+1}\). If \(R^c\) is replaced, let
\(e_R\) be its order as computed above, otherwise set \(e_R=0\).
The total extra order is
\[
 e=e_R+\sum_{\rm replaced}e_l>0.
\]

Discard terms with \(e\ge b\). For the others we recover all
\(R_0^c\) by spending at most \(e\) of the \(b\) spare
determinants at each \(j\). Put the unused excess above
\(b-e\) into the tangential coefficient.

\proofstep{The remaining normal order after correction.}
The newly exposed normal multiplicity \(m_*\) satisfies
\[
 m_*\le e_R+\sum_{\rm replaced}(d_l+e_l)
       \le (2Q+1)e .
\]
At each normal variable use one copy of the corresponding
\(T_s\) and the preceding coordinate identity. This replaces it
with a sum of order-one leading expressions divided by \(w\)
and multiplied by tangential coefficients. There are enough
copies at every position to leave \((2Q+3)(b-e)\) copies, with
the surplus also tangential. After absorbing \(w^{m_*}\) into the order-one leading
expressions, the external power is
\(i'=i+\sum_{\rm replaced}d_l+e-m_*\ge0\).

Adding the degrees of the untouched factors and the
\(m_*\) order-one factors yields precisely
\[
 i'+\sum_{\rm untouched} d_l+m_*=L+e .
\]
The residual determinant reserves now have parameter \(b-e\),
so each error term is in the induction at the smaller parameter. Induction lifts them
to precision \((L+e)+(b-e)=L+b\) at the chosen plane and with
at least that order at competitors. Subtract those lifts from
the raw product. Finally one can project onto the prescribed
multidegree, since every substitution is linear within each
true letter and preserves its total degree.  Its use requires generating symbols
with the displayed determinant powers in the available degrees.

\end{proof}

\paragraph{Extracting determinants from a prefix.}

The correction must also preserve the orders at the cuts of a path. Extracting powers of the determinant from a prefix measures how much of its order remains.

\begin{lemma}[Determinant divisibility from a prefix order]\label{lem:plane-prefix-determinant}
In the plane variables
before a cut, let \(A\) be a list of \(r\) distinct vector
letters, let \(Z\) be the square matrix of their rows, and
let \(B\) be the greatest total degree in all other vector
letters before the cut in any monomial of an integral polynomial
\(P\). Coefficients may themselves be polynomials in auxiliary
variables over the DVR. Suppose that under a generic conjugate
of \(\operatorname{diag}(1,\ldots,1,T)\) acting on the prefix
vectors and contragrediently on its covectors every Laurent
power in \(P\) is at least \(\theta>B\). Then
\[
                    P=(\det Z)^{\theta-B}P_1
\]
integrally. Dividing by this factor lowers the minimum Laurent
power on any other subset of letters by at most \(\theta-B\).
\end{lemma}

\begin{proof}
\proofstep{The determinant identity on the prefix chart.}
For the first assertion localize the polynomial ring at the
prime of the generic determinant \(\pi=\det Z\). This ring
is a DVR. Indeed the generic determinant is irreducible over
the fraction field of the coefficients and primitive over the
coefficient UFD. Irreducibility follows by expanding in the
last row, whose complementary minors have no common factor.
For \(r>1\) two distinct such minors are nonassociate irreducibles
by induction, and the case \(r=1\) is immediate. Only the finitely
many variables occurring in the calculation are needed.

In the
residue field of this localization \(Z\) has rank \(r-1\)
since a nonzero minor of size \(r-1\) is a unit there.
Choose \(g\in\mathrm{GL}_r\) over the localization whose last
column spans its kernel modulo \(\pi\), by lifting a basis of
the residue space. Write
\[
 H_T=g\operatorname{diag} (1,\ldots,1,T)g^{-1}.
\]
Specialize the generic conjugating
matrix to this \(g\). All coefficients remain integral and the
lower Laurent-power bound persists, as the only matrix denominator
is a power of the unit \(\det g\).

In the transformed polynomial evaluate the input prefix rows at
\(xH_\pi^{-1}\) and the input covector columns at \(H_\pi y\).
At \(T=\pi\) the result is exactly the original \(P\).
All covectors are integral. Each row of \(Zg\) has last entry
divisible by \(\pi\), so the \(r\) chosen input rows are
integral as well. Each remaining prefix vector can cost one
power of \(\pi\) per coordinate factor. Their combined
degree is still at most \(B\) under a linear change within
each letter. For a term \(A_l\) of Laurent degree \(l\ge\theta\), the possible denominator loss is therefore at most \(B\), giving
\[
 \operatorname{ord}_\pi(A_l|_{T=\pi})\ge l-B\ge\theta-B.
\]

This gives the claimed divisibility in the localization and
hence in the original UFD. For the last assertion apply the
generic conjugate to any subset of the determinant rows. Their
wedge is transformed by an exterior power of \(H_T\), in which
the largest Laurent exponent is at most one. The determinant
has minimum exponent at most one there. Minimum exponents
add on nonzero products over a domain, giving the asserted
bound for \(P_1\). The case \(P=0\) uses \(P_1=0\).
The other cuts are controlled simultaneously.

\end{proof}

\begin{lemma}[Allocating the determinant reserves]\label{lem:plane-determinant-reserve}
Assume the signed-prefix order and the one-sided coarse degree
bound of \eqref{eq:plane-prefix-degree-assumptions}, with
\(p\ge13\), \(k\ge3\), \(Y\ge2\), and the bounds on size and precision in Section~\ref{sec:host-edge-tests}.
The residue-tangential product has a factor
\(F_*\prod_{s\ {\rm top}}T_s^{kQ}\) to precision \(p^{c'}\)
after substituting the original radices. Here
\[
 F_*=R_0^{c'}\prod_jD_j^{c'}
                         \prod_{s\ {\rm top}}T_s^{(2Q+3)c'} .
\]
These factors are available in the stated degrees for
Lemma~\ref{lem:plane-normal-correction}.
\end{lemma}

\begin{proof}
\proofstep{Allocating the determinant factors.}
For an application at the coarse degrees, let \(z\) be the last
position of the outgoing frame in a segment. At a top \(s\)
after the incoming frame take \(A\) to be the preceding \(r\)
even tops and \(u\) their earliest position. Write
\(B_o=\deg(u)/p^{z-u}\), with the true top degrees between
\(a_0\) and \(2a_0\), where \(a_0\) is the lower bound for
the plane model (at the host it is \(a-c_G\), up to integer
rounding). Put \(Y=Y_2\). Let the minimum of the Laurent exponents at the
prefix strictly before \(s\) be \(\operatorname{ord}_{<s}\).
The conditions for this prefix and the coarse degrees on the
incoming frame and intervening tops are
\begin{equation}\label{eq:plane-prefix-degree-assumptions}
 \operatorname{ord}_{<s}\ge\theta
      =\left\lceil B_o/2-\left\lfloor B_o/(Yk^2)\right\rfloor\right\rceil,
 \qquad
 d_y\le1.01\deg(y)/p^{z-y}.
\end{equation}
The order bound holds on sums because it excludes every lower
Laurent power.

In particular the stated one-percent two-sided hypothesis suffices.
Every earlier even top outside \(A\) belongs to these lists,
because the outgoing frame ends the segment and has only
\(r\) members. Such tops have \(y\le u-2\), so
\[
 B\le \sum_{y<u,\ y\ {\rm even}} d_y
       \le 2.02 B_o\sum_{j\ge1}p^{-2j}
       =2.02 B_o/(p^2-1).
\]
For \(p\ge13,\ k\ge3,\ Y\ge2\), this leaves
\(\theta-B> 2B_o/5\). In particular these are large
reserves even for the earliest row, since \(B_o\ge a_0/Q\).
For example let \(c'=8(D_{\rm II}+1)+rN c_0+
2rn'\lfloor c_G\rfloor\) and set
\[
 m_*=(N_1+1)\big((2Q+3)c'+kQ+c'\max_j\alpha_j+3c'\big).
\]
As \(\max\alpha_j\le1+rkQ\), \(N_1^2\le rN\), and
\(k^3r^2NQ^5 c'Y=o(a_0)\) by
Section~\ref{sec:host-edge-tests}, we have \(N_1m_*=o(a_0/Q)\).

At each such cut extract \(m_1=(2Q+3)c'+kQ+c'\) copies of the
ordinary determinant on the chosen rows. From the incoming frame
extract \(m_2=c'\max_j\alpha_j+2c'\) copies of its determinant.
This frame consists of the preceding \(r\) rows at the first cut.
The total number extracted in a segment is at most
\[
 N_1m_1+m_2\le m_*.
\]
Each copy costs at most one in every
prefix minimum and cannot increase the degree in any subset of
the other vector variables, since highest degrees in a subset
add on nonzero products. The absolute lower bound
\(\theta-B>2B_o/5\) permits these successive extractions.

The
original plane-versus-complement inequalities have margin
\(m_*\), which absorbs their loss from the extractions at each cut,
including the cuts that end at a covector.
The determinants continue to give the required tangential
factors after radix substitution, to the precision needed for
the symbol. Indeed put \(P_s=p^{z-s}\). The image of a
determinant on the \(r\) preceding rows is \(\det(U_s^{[P_s]})\).
Over \(O\) this is
\[
 T_s^{P_s}+pE
\]
with \(E\) integral,
since Frobenius commutes with determinants in characteristic
\(p\). In a power of \(h+c'\) such factors, terms using more
than \(c'\) copies of \(pE\) vanish modulo \(w^{c'}\).
The remaining terms have the common factor \(T_s^h\), with
integral tangential quotient.

Take \(h=(2Q+3)c'+kQ\).
The coarse determinant of the incoming frame has image
\(D_{j-1}^{p^{z-z_{j-1}}}+pE'\). The same argument gives a
factor \(D_{j-1}^{(1+\alpha_{j-1})c'}\) from the \(m_2\) copies.
Every top occurs after the incoming frame in exactly one segment
per cycle. Around the cycle the factored expression is therefore
\[
 F_*\prod_{s\ {\rm top}}T_s^{kQ}\,H,\qquad
 F_* = R_0^{c'}\prod_jD_j^{c'}
                 \prod_{s\ {\rm top}} T_s^{(2Q+3)c'},
\]
for an integral tangential polynomial \(H\), and the discrepancy
has coefficients in \(p^{c'}O\).  The factor \(\prod_s T_s^{kQ}\) will be used
on the digit labels.

\end{proof}

\paragraph{Integral generators after determinant extraction.}
\hypertarget{source-plane-symbols}{}
This part of the proof of Proposition~\ref{prop:residue-plane-neighbor-test}
gives the integral source polynomial on a single plane.
The determinant factors extracted above remain part of the
symbol. The result of the lifting is as follows.

\begin{lemma}[Lifting a plane symbol]\label{lem:lift-plane-symbol}
Fix a selected plane and the determinant factors of
Lemma~\ref{lem:plane-determinant-reserve}. Suppose that a residue symbol
\(\sigma\) has normal degree \(o\) and belongs to the prescribed
span of digit monomials times products of segment symbols,
with their prefix and multidegree conditions. Write \(\Pi\) for the selected plane. At precision \(c'\) there is one integral polynomial \(P\) with the following two restrictions
\[
 P|_\Pi\equiv w^o\sigma\pmod{w^{o+c'}},\qquad P|_{\Pi'}\equiv0\pmod{w^{o+c'}}\quad(\Pi'\text{ competing}).
\]
 It may be projected to the required
true multidegree. The column factors used in the lift have
prescribed order at most \(2Q\).
\end{lemma}

\begin{proof}
\proofstep{A paired path after the determinant extraction.}
Split the coordinate module into its plane and complement,
\[
 O^k=E\oplus C,\qquad E=O^r,\quad C=O^K,\quad K=k-r.
\]
At a top indexed by \(s\)
the even letters are vectors and the odd letters are covectors.
Fix one segment with end index \(z\). Its prescribed coarse
degree in a letter is the sum of its \(E\)-degree and its
\(C\)-degree. A segment symbol belongs to the span of the products \(A(E)G(C)\) satisfying the prefix bounds. To put the summands at the prescribed coarse
degrees, decompose \(A\) and \(G\) into homogeneous components
in every letter and keep the pairs whose degrees sum to the
prescribed ones. If \(\pi_\alpha\) extracts a multihomogeneous component, each prefix substitution \(g\) is linear within its letters. Hence
\[
 \pi_\alpha(gF)=g\,\pi_\alpha(F).
\]
Each component retains the same bounds. The
normal-degree cutoff is a condition on the original segment
symbol used in the rank argument. The realization here
applies to this whole span after the multidegree projection.

Apply the determinant extraction to \(A\) in each summand.
The determinants have fixed degrees in each coarse letter, so
the quotient \(A'\) also has fixed degrees. For every top \(s\)
after the incoming frame, let \(a_{<s},a_{\le s}\) denote the
minimum Laurent exponents of \(A'\) under the two prefix actions
ending before and at \(s\). Let \(g_{<s},g_{\le s}\) denote the
corresponding maximum exponents of \(G\). The remaining
separation says
\[
  g_{\le s}\le a_{<s}\quad(s\ {\rm even}),\qquad
  g_{<s}\le a_{\le s}\quad(s\ {\rm odd}).
\]
The complemented factor has degree zero in the incoming frame.

\proofstep{Integral basis after removing the determinant factors.}
Use the integral Pieri basis from
Section~\ref{sec:column-lifts} for the two factors separately.
The prefix filtration used there also works with a lower
bound on the last entry of the intermediate weight. In each
column the signed subset weight on one coordinate lies
between the contributions from the smallest and largest
entries of the intermediate weight. In the product it lies
between those two entries. This gives both bounds upon
conjugating the one-coordinate torus.

Conversely, over the
fraction field, the requirement that every conjugate torus
act within the specified range selects exactly the isotypic
summands whose first and last entries lie within that range.
Indeed the permitted subspace is invariant under the prefix
group. On each copy of \(V_\lambda\) in \(V_\lambda\otimes M_\lambda\), the extrema of the coordinate weights \(\mu\) are
\[
 \max\mu=\lambda_1,\qquad \min\mu=\begin{cases}\lambda_r,&\text{plane},\\\lambda_K,&\text{complement}.\end{cases}
\]
Complete reducibility gives the assertion. Splitting
successively by Pieri at the prefixes as in that section proves
the dimension formula for the intersection. The allowed basis
words span the intersection over \(O\) because the full basis is
unimodular.

Use the actual lower prefix exponents of \(A'\) as its lower
bounds and the actual upper prefix exponents of \(G\) as its
upper bounds, discarding zero summands. If \(\nu\) and
\(\kappa\) are the paths in any resulting product of basis
words, then for even \(s\) one has
\[
 \kappa_1(\le s)\le\nu_r(<s).
\]
For odd \(s\) one has
\(\kappa_1(<s)\le\nu_r(\le s)\).
These inequalities give exactly the interlacing across the
boundary between the \(r\)-th and \((r+1)\)-st rows of the
concatenated path \(\lambda=(\nu,\kappa)\). Each separate path already interlaces within its own rows.

For the boundary, if \(\lambda_i(s)\) denotes the combined
row weight after \(s\), the inequalities say
\[
 \lambda_r(s-1)\ge\lambda_{r+1}(s)\quad(s\ \mathrm{even}),
 \qquad
 \lambda_r(s)\ge\lambda_{r+1}(s-1)\quad(s\ \mathrm{odd}).
\]
The two row weights are nondecreasing at vector steps and
nonincreasing at covector steps. Thus the displayed boundary
inequality and the separate interlacings give dominance
both before and after each step. In the incoming frame
the complement stays at zero and plane entries are nonnegative
and increasing, so interlacing also holds there. Each strip's
size is the sum of the two sizes, and the concatenated path
starts at zero.

\paragraph{Full order of the two-block columns.}
\label{par:plane-column-orders}
Fix a column \(b\) of a concatenated path, and write \(h\) for
its height at a cut. A positive column uses the exterior
degree \(h\); a nonpositive column uses dual exterior degree
\(k-h\), as in Section~\ref{sec:column-lifts}. All changes of
height are by one, rising at a vector and falling at a
covector. For a subset \(I\) of size \(h\) in a positive column,
\[
 \#(I\cap C)=h-\#(I\cap E)\ge h-\min(r,h).
\]
Here the intersections count coordinate indices in the two summands. Equality characterizes the subsets of smallest normal cost. For a
nonpositive column
\[
 h>r
\]
at every cut. Indeed it starts at
\(k\) and loses at most one per letter, and the segment length
plus \(r\) is at most \(K\). The subsets of smallest cost
in the dual tensor use only complement indices.

The costly transitions have precisely the two forms
\[
 h\longrightarrow h+1\quad(h\ge r,\ \text{vector}),\qquad h\longrightarrow h-1\quad(h\le r,\ \text{covector}).
\]
In a nonpositive
column only the first possibility occurs. This is the normal
coordinate selected when both exterior subsets have their
prescribed plane occupancy. For a positive column, define the
defect of a subset as \(\min(r,h)\) minus its plane occupancy.
For a nonpositive column its defect is its plane occupancy.
In either case, along a sequence of wedge and contraction
terms, write \(n_i=1\) if the \(i\)-th letter
uses a normal coordinate and \(c_i=1\) if that step is
costly, taking both indicators to be zero in the other
case. The defect \(\delta_i\) then satisfies
\[
                    \delta_i-\delta_{i-1}=n_i-c_i .
\]
This follows from the change in occupancy in the wedge
or contraction operation and the change in the target
occupancy \(\min(r,h)\).

For example, at a positive rise from \(h<r\), adding a
complement index increases the defect by one. At a positive
fall from \(h\le r\), removing a complement index decreases
it by one. When \(h>r\) a rise on a plane index decreases the
defect and a fall on that index increases it. The dual
operations wedge on a covector and contract on a vector,
which increase and decrease its plane occupancy when the
plane index is used.

\proofstep{The allowed path of the restored column.}
The segment letters are the increasing enumeration of a set of
distinct top positions \(s_i\). Their radices
\[
 \omega_i=p^{z-s_i}>0
\]
are strictly decreasing. Let
\(\delta_i\) be the defect just after letter
\(i\), with \(\delta_0=0\). The excess of weighted normal cost
over the weighted count of costly steps is
\[
 \sum_{i<M}\delta_i(\omega_i-\omega_{i+1})+\delta_M\omega_M
\]
for a segment of \(M\) letters. Thus the lowest part is
computed by keeping defect zero throughout. In positive
columns the plane operations then take place with no
complement indices in the subset. All remaining operations
act on the complement past the plane indices.

In the
nonpositive columns the plane exterior degree in the dual
tensor stays zero. With the right wedge and contraction
conventions, an operation on a complement index counts only
larger complement indices for its sign, and every active
operation on a plane index counts only larger plane indices.
The defect-zero recursion therefore splits without a sign
change into the two column recursions for \(\nu\) and
\(\kappa\). For terminal coordinates take their union in the
corresponding blocks. Writing \(\operatorname{in}\) for the component of smallest normal degree, we obtain for the column polynomials \(C_b\)
\[
 \operatorname{in}\Bigl(\prod_b C_b\Bigr)=\prod_b\operatorname{in}(C_b).
\]
The defect-zero recursion identifies the right side with the prescribed paired word.

For the true lift, use the same exterior recursion in the
designated twisted frame. Between successive used indices
\(s<t\), raise every entry to \(p^{t-s}\), then perform the
operation at \(t\). After the last operation power up to \(z\).
Take \(\omega_i=p^{z-s_i}\), where \(s_i\) is the index of
letter \(i\). The column has degree \(\omega_i\) in each used
true letter. Its prescribed order is
\[
             d_b=\sum_{i\ {\rm costly}}\omega_i\le 2Q .
\]
The bound follows from the strict spacing of the used integer
indices and the geometric sum starting at the largest power,
which is at most \(Q\).

\proofstep{The order in a neighboring plane.}
For any tested plane, measure the running exterior tensor in
its own twisted frame. Suppose at a stage the common order is
at least \(d\), and every entry with positive defect has order
at least \(d+1\). This holds before the first operation.
For the passage \(s<t\) put \(P=p^{t-s}\) and write the
adapted coordinates at \(s\) as \(w^d U\). If \(L_s,L_t\)
are the changes of exterior coordinates from the designated
frame to the tested frame at the two times, then
\[
       (w^d U L_s^{-1})^{[P]} L_t
                 = w^{Pd}\big(U^{[P]}+pV\big)
\]
with \(V\) integral. Entrywise Frobenius is additive modulo
\(p\), and \((L_s^{-1})^{[P]}\equiv L_t^{-1}\pmod p\)
because the matrices and their inverse minors come from the
unramified ring and its successive Frobenius twists.

Since
\(p=w^2\) and \(P\ge2\), the terms in positive-defect entries
now have order at least \(Pd+2\). At a free step the free
coordinates take defect-zero subsets to defect-zero subsets,
and the normal coordinates contribute a factor \(w\).
At a costly step the free part kills the entries at defect
zero. The normal part carries them to the new defect-zero
subsets and contributes \(w\). Use as a common lower bound immediately before the linear step the order \(Pd\).
After a costly step the required order is \(d'=Pd+1\).
Thus the guaranteed lower bounds for the two contributions are
\[
 \begin{aligned}
 \text{defect zero, normal factor }w &: \quad Pd+1=d',\\
 \text{positive defect, integral factor} &: \quad Pd+2\ge d'+1 .
 \end{aligned}
\]
Only the first contribution can affect the leading coefficient.
At a free step the required order stays \(Pd\), and a normal
factor again raises it. These inequalities apply in every frame.

At the designated frame the change-of-frame matrices are
identities. After division by the prescribed power of \(w\)
the leading recursion simply keeps the defect-zero part. The
powering between steps and up to \(z\) commutes with its
additions and multiplications modulo \(w\), where the
characteristic is \(p\). Its result is therefore the
defect-zero coarse column with the substitution
\[
 X_{i,a}\mapsto X_{s_i,a}^{p^{z-s_i}}.
\]

Consequently the integral true column is a factor of order
\(d_b\) for the correction lemma. The product of its
normalized leading expressions is the desired paired symbol
after this substitution, with the full normal degree equal
to the sum of the column orders.  \label{cor:plane-symbols-1}

\paragraph{Digits, normal degree, and the final projection.}
Let \(X^\ell\) be the common digit monomial. At every top
each digit is less than the letter's radix, which is at
most \(Q\). Thus its total normal degree at that top is at
most \(kQ\). Spend one \(T_s\) for each normal digit factor
and use the leading order-one formulas for that coordinate.
If \(d_s\) normal digit factors are used at this frame, the reserve of \(T_s\) becomes
\[
 T_s^{kQ-d_s},\qquad kQ-d_s\ge0.
\]
It and the tangential digit variables are integral tangential coefficients. This
expresses \(X^\ell\prod_s T_s^{kQ}\) as a sum of products of
normalized leading expressions of universal order-one
factors with integral tangential coefficients. Every term has
the full normal degree of the digit monomial.

\proofstep{Recovering the true letter degree.}
Now start with a symbol \(\sigma\) of normal degree \(o\),
presented as \(X^\ell\) times a sum of products of the
transformed segment symbols. The extracted determinant
product is common to these summands. In the distinguished
coordinates it gives the identity
\[
  \sigma =
  F_*\,H\,X^\ell\!\prod_s T_s^{kQ}\,S + p^{c'}Y',
\]
where \(S\) is the sum of products of transformed residual
segments and \(Y'\) is integral. After expanding residual
segments in paired words, every summand on the right outside
the error has normal degree equal to the sum of the
appropriate column and digit orders. Since \(F_*,H\) are
tangential, take the component of normal degree \(o\).
The resulting main terms, after multiplication by \(w^o\),
are products of leading expressions of combined order \(o\),
with the reserves for precision \(c'\) contained in \(F_*\).
The error acquires order at least \(o+2c'\).

Apply the correction lemma to each surviving term, then add.
For the resulting integral true polynomial \(P\), the two tests read
\[
 \begin{aligned}P|_{\rm designated}&\equiv w^o\sigma\pmod{w^{o+c'}},\\ P|_{\rm competing}&\equiv0\pmod{w^{o+c'}}\qquad(\text{each competitor}).\end{aligned}
\]
If \(\sigma\) has a prescribed true multidegree,
project the resulting polynomial on that multidegree.
Every plane substitution is linear within each true
letter, so the projection commutes with all the substitutions
and preserves all divisibilities.  Its uniform order
bound on a key factor is \(2Q\), and its precision is the
prescribed \(c'\) at every plane. The determinant estimate above
also proves the one-sided version used for the dual test in
Section~\ref{par:plane-reassembly}. \label{cor:plane-symbols-2}
\end{proof}

\begin{paramdefinition}[Frames, coarsening and sizes for the plane lattice]\label{hyp:coarsened-plane-frames}\label{par:plane-rank}
\hypertarget{source-plane-rank}{}
In the plane construction let \(r\) be the plane dimension and \(K\) its codimension in dimension \(k=r+K\). Write \(a\) for the lower bound on the boundary degrees \(b_x\).
Let \(J\) be the number of frames, \(N\) the number of tops in a
cycle, and \(N_s\) an upper bound for the number of letters in a
segment. In the count use the lifted frame endpoints \(s_j,z_j\)
with \(s_j\) the first position of frame \(j\). Segment \(j\)
starts at \(s_{j-1}\) and ends at \(z_j\). Write \(Q\) for the
original bound on \(p^{z_j-x}\). The coarsening uses
\[
 A=\frac{a}{8Yk^2N_s},\qquad L=\lfloor\log_p A\rfloor,\qquad
 P_{j,x}=p^{L+s_{j-1}-x},\qquad H_j=p^{s_j-s_{j-1}}.
\]
Indeed the extra power at segment \(j\) is the largest power of
\(p\) at most \(\lfloor A/p^{z_j-s_{j-1}}\rfloor\).

Since \(A\ge Q\),
this equals \(p^{L-(z_j-s_{j-1})}\). In particular \(P_{j,x}\le A\),
\(P_{j+1,x}=P_{j,x}H_j\) on frame \(j\), and \(\prod_j H_j=q\).
Use the separation bound \(p^{L_1-1}\ge128 k^2N_sY^2\) and the
size bound
\begin{equation}\label{eq:plane-count-size}
  a\ge32pYk^2N_sQ^2(E_0+m_*+1).
\end{equation}
These are the separation and size inequalities for the plane model.
The argument below uses them with \(p\ge13\), \(r\ge2\),
\(Y\ge1\), \(N_s+r\le K\), \(a\le b_x\le 2a\), and the lifted
frame assumptions of Section~\ref{sec:plane-lifts}.
Put \(D_k(b)=\binom{b+k-1}{k-1}\).
\end{paramdefinition}
For the selected cycle, the separation and size comparisons in the proof of Proposition~\ref{prop:host-small-prime-test} use \(Y=x^{26}\), \(p^{L_1-1}\ge256k^2NY^2\) and \(\log Q=o(\log(a/b'))\). They hold uniformly over the admissible frames.


\begin{lemma}[Frame and digit degrees in the coarsened lattice]\label{lem:plane-digit-frame-lattice}
For the coarsened lattice of Definition~\ref{hyp:coarsened-plane-frames},
the late degrees satisfy
\[
 a'_x=b_x-E'_x-E_x>0.
\]

The digit and outgoing choices produce the same total degree
\(a'_x\) in the true coordinate at each top, with incoming
degree \(b_{j+1,x}\ge W_{j+1,x}\ge kN_s+1\).
The corresponding segment word families are independent over
the residue field in the fixed digit and degree choices.
\end{lemma}

\begin{proof}
\proofstep{The coarse degrees at the later tops.}
At each late top \(x\) of segment \(j\), meaning one after the
incoming frame, form the ordinary plane determinant of the last
\(r\) preceding even rows and give it the exponent
\[
 g_{j,x}=\left\lfloor\frac{b_u}{k^2Yp^{z_j-u}}\right\rfloor,
 \qquad u=\text{the earliest of those rows}.
\]
Write \(\mathcal D\) for the product over all these determinants
with their original radices \(p^{z_j-y}\).

At a top \(y\) its
true degree \(E_y\) is at most \(4N_s a/(k^2Y)\). A determinant
contributes at most \(g_{j,x}p^{z_j-y}\le 2a/(k^2Y)\), since
\(y\ge u\). Each top occurs in at most two segments. Use coarsened
true degrees \(a'_x=b_x-E'_x-E_x\), where \(E'_x\le E_0\) is the
prior cleaning budget. In particular
\[
 a'_x\ge a/2,
\]
since
\(4N_s/k^2\le 1/4\) follows from \(N_s+2r\le k\),
\(r\ge2\), and \(E_0\le a/32\).

At each top there is a unique occurrence after the incoming
frame. Choose a coordinate digit vector \(\ell_x\in[0,P_{j,x})^k\)
at this occurrence, with total \(|\ell_x|\equiv a'_x\pmod{P_{j,x}}\).
All digit choices have \(|\ell_x|<kP_{j,x}\le a'_x\). Write
\[
 n_x=(a'_x-|\ell_x|)/P_{j,x}.
\]
At a top strictly between frames
the degree in the segment is \(n_x\).

At an outgoing frame row
of segment \(j\) choose a degree \(e_x\in[0,rH_j)\) with
\(e_x\equiv n_x\pmod{H_j}\). Use degree \(b_{j+1,x}=(n_x-e_x)/H_j\)
when the same top reappears in the incoming frame. In fact
\[
 b_{j+1,x}>\frac{a'_x}{P_{j+1,x}}-k-r,
 \qquad
 W_{j+1,x}=\left\lfloor\frac{a'_x}{P_{j+1,x}}\right\rfloor-k-r-1,
 \qquad b_{j+1,x}\ge W_{j+1,x}.
\]
The strict inequality uses \(H_j\ge1\) and \(e_x<rH_j\). Our size
bounds give \(a'_x/P_{j+1,x}\ge4Yk^2N_s\), so
\[
 W_{j+1,x}\ge m_0=kN_s+1
\]

for \(m_0=kN_s+1\). The digit and frame degrees thus add to exactly
\(a'_x\) at every top, with nonnegative incoming padding.

\paragraph{The segment words.}
Fix the digits of the intervening tops in segment \(j\). The pre-outgoing path degree at \(x\) is
\[
 B_x=\begin{cases}n_x,&\text{intervening top},\\ W_{j,x}-m_0,&\text{incoming row}.\end{cases}
\]
Write \(S=\sum_x B_x\)
for this sum before the outgoing frame. Consider paths \(\mu\) in
\(k\) rows with these degrees. The last \(r\) steps are vector
steps of variable size. Just after the \(i\)-th outgoing step
require that the first \(i\) entries are all \(S\), with indices
starting at one.

Write \(\nu,\kappa\) for the restrictions to the
first \(r\) and the last \(K\) entries. Each is a signed Pieri
path with sizes determined by its weight changes. Each covector step creates at most one further negative entry. Thus
\[
 \#\{i:\mu_i<0\}\le N_s\le K-r,
\]
so the first \(r\) entries of the dominant weight \(\mu\) stay nonnegative. They stay at most \(S\) before the outgoing
frame by the sum of previous vector degrees. They stay at most
\(S\) afterward by the head condition and dominance. The
complement has zero degrees in the incoming frame because the
first \(r\) steps are vectors.

Impose the admissible path inequalities. At every
late top \(x\), \(\nu_r\) before the step is at least \(2B_u/3\)
for the earliest \(u\) of the preceding \(r\) vectors. The
complement increase at a vector and the plane decrease at a
covector are at most
\[
 \frac{\mathcal L a}{pkP_{j,x}}+
   \mathbf 1_{x\text{ even}}\frac{4a}{P_{j,u}},
 \qquad \mathcal L=C_{\mathcal L} N_s^2\log a.
\]
Let \(\mathcal T_j\) be the set of such paths with a choice of
standard terminal tableau for \(\kappa\). Its size depends only on
the intervening degrees and the fixed data.

For the plane word use
the tableau of its terminal rectangle \((S,\ldots,S)\), which is
unique. Let \(d_i\) be the full strip size at outgoing row \(i\).
It is at most \(2S\). At a vector step from \(\mu\) to \(\mu'\),
interlacing gives \(\sum_h(\mu'_h-\mu_h)\le\mu'_1-\mu_k\).
Here \(\mu'_1=S\), and \(\mu_k\ge-S\) since the total of all
covector degrees is at most \(S\).

\proofstep{Counting the outgoing decorations.}
For each member of \(\mathcal T_j\), choose outgoing decoration
exponents \(v_{i,h}\ge0\) for the \(r\) plane coordinates of
outgoing row \(i\), with
\begin{equation}\label{eq:plane-outgoing-decoration}
 v_{i,h}+2S<H_j,\qquad d_i+\sum_h v_{i,h}=e_i.
\end{equation}
Use the product of the plane and complement words, the outgoing
monomial with these exponents, and the same incoming factor for
all paths at the chosen degrees. The incoming factor is the
plane determinant of the incoming frame raised to \(m_0\), times
one coordinate of each incoming row to degree \(b_{j,x}-W_{j,x}\).
The product is in the indicated lattice \(M_j\). The signed strip
interlacing at the plane boundary supplies the two inequalities
between plane and complement at each late cut. The incoming determinant power shifts every plane coordinate of the weight by \(m_0\).

If the old minimum is \(\nu_r\), then
\[
 \min_i(\nu_i+m_0)=\nu_r+m_0.
\]
The other added vector factors can only increase this minimum.
The path bound on \(\nu_r\) before the step gives the absolute
lower cutoff independently. The normal degree of the product
at a vector is the complement strip size and at a covector it
is the plane strip size. The path inequalities give the normal
cutoff. Finally the exponents in each outgoing plane coordinate
are less than \(H_j\) by \eqref{eq:plane-outgoing-decoration} and
the bound
\[
 d_i\le2S.
\]


\proofstep{Independence in fixed degrees.}
These products are independent at each fixed set of degrees,
also over the residue field. First omit their common nonzero
incoming factor. In the \(i\)-th outgoing plane row specialize
all but coordinate \(r+1-i\) to zero. For each column
\(c\in[1,S]\) of the plane word, the height increases by zero
or one at each outgoing step. A rise wedges the sole surviving
coordinate there. Reading back from the final full exterior
subset, remove each of the coordinates at its rise. This leaves
precisely the terminal subset for the column at the pre-outgoing
cut, the complement of its rise coordinates. The coefficient is
\(1\) or \(-1\) by the wedge rule.

If \(R_c\) is the set of outgoing rises in column \(c\),
its pre-outgoing terminal subset is
\[
       [1,r]\setminus\{r+1-i:i\in R_c\}.
\]
Thus the specialization of the product word is, up to sign,
\[
 \prod_{i=1}^r x_{i,r+1-i}^{\,\#\{c:i\in R_c\}}
\]
times the
truncated plane word with those terminal subsets and the whole
complement word. The variables \(x_{i,r+1-i}\) are the surviving
outgoing coordinates.
The terminal subsets of the truncated plane word are standard.
At every outgoing cut the height on the left is at least the
height on the right.

If their current heights are \(h_L\ge h_R\), the remaining rises to the common terminal height obey
\[
 r-h_L\le r-h_R.
\]
Let \(D_L,D_R\) be the sets of removed coordinates in two adjacent columns. The surviving coordinates run in reverse order, so for every initial interval \([1,b]\), 
\[
 |D_L\cap[1,b]|\le|D_R\cap[1,b]|.
\]
 The retained
subsets therefore have weakly increasing entries across a row.
The truncated plane path and its subsets recover all
plane heights, since each missing coordinate specifies the time
of one rise. The two paths and their tableaux are therefore
recovered from the standard words in the variables other than
the outgoing plane variables.

To include the decorations in this argument, take the terms of
least total degree in the discarded coordinates in a putative
relation. For each tuple of exponents in those coordinates and
in the surviving outgoing plane coordinates, the displayed
specializations are independent by the standard word basis
in the remaining variables. Once the word indices are known, let \(a_i\) be their fixed antidiagonal exponents. The observed exponents \(b_i\) then recover each decoration by
\[
 d_i=b_i-a_i.
\]
This proves
independence and also explains why it holds with the same
coefficient ring as the residue rank calculation.

\end{proof}

\paragraph{Weight of paths at the segment cutoffs.}
\label{par:plane-typical}
\begin{lemma}[Weight of the paths obeying the cutoffs]
\label{lem:good-plane-path-count}
For segment \(j\) of Definition~\ref{hyp:coarsened-plane-frames}, let \(\mathcal T_j\) be the paths satisfying the degree cutoffs, with a choice of standard terminal tableau for the complement weight \(\kappa\). With \(D_k(b)=\binom{b+k-1}{k-1}\), their number satisfies
\begin{equation}\label{eq:plane-good-path-count}
 |\mathcal T_j|\ \ge\ (1-a^{-N_s^2})
    \prod_{x\text{ between frames}}D_k(n_x)
    \prod_{x\text{ incoming}}D_k(W_{j,x}-m_0).
\end{equation}
\end{lemma}

\begin{proof}
We prove this at the simultaneous cutoffs above. Let \(m\le N_s\)
be the number of letters in the segment and number the first
\(m-r\) steps before the outgoing frame. Use the hypotheses
\(k,\mathcal L\ge1088N_s^2\log a\). From the degree-loss bound
\(E'_x+E_x\le5b_x/(kY)\) we have \(a'_x\ge9b_x/10\).
For either kind of pre-outgoing row,
\[
 \frac{4a}{5P_{j,x}}\le B_x\le\frac{2a}{P_{j,x}},\qquad B_x\ge k.
\]
In fact the floor and padding lose at most
\[
 4kN_s\le b_x/(10P_{j,x}),
\]

since \(b_x/P_{j,x}\ge8Yk^2N_s\). Thus degrees of successive
rows of the same parity grow by at least
\[
 \rho=p^2/5>13.
\]

For any letter \(x\), each preceding pre-outgoing degree is at
most \(4\beta_x/p\), where \(\beta_x=a/P_{j,x}\), by strict
increase of the top positions.

First count all paths with the outgoing head conditions,
weighted by the number of tableaux of the terminal complement.
Before the outgoing frame the signed Pieri rules decompose the
tensor product of the indicated symmetric powers of the standard
module and its dual. Each path to a weight \(\lambda\) has weight
\(\dim V_\lambda\). After \(t\) outgoing steps remove the first \(t\)
entries, all equal to \(S\), and call the remaining row list the
tail.

At the next step its old and new tails \(\eta,\eta'\) obey
\[
 \eta_i\ge\eta'_i\ge\eta_{i+1}.
\]
Conversely every such
interlacing tail reinstates a vector strip on inserting the
extra head equal to \(S\), because all entries were at most
\(S\). Thus at the outgoing steps the ordinary branching rule
from \(\mathrm{GL}_{k-t}\) to \(\mathrm{GL}_{k-t-1}\) preserves the
total weight. These multiplicity-one rules and the Weyl dimension
formula apply to rational weights by a determinant twist and
duality. \cite{fulton-harris}
The last tail is the complement, whose tableau count is its
Weyl dimension. The total count before imposing the good path
inequalities is exactly
\[
 R=\prod_{x\text{ pre-outgoing}}D_k(B_x).
\]


\proofstep{The constraints on a path and its two signs.}
For a given path and sign \(\epsilon\in\{+,-\}\), let
\(\alpha^\epsilon\) be the nonincreasing list of positive parts
or of absolute negative parts of its tail. Up to the outgoing
frame the tail is the full list of \(k\) rows. Let \(\mathcal B^\epsilon\)
be the list of previously used pre-outgoing degrees of the
corresponding sign, with vectors positive and covectors negative.
For every \(T\ge0\) put
\[
 \delta_\epsilon(T)=\sum_{B\in\mathcal B^\epsilon}\min(B,T)
                    -\sum_h\min(\alpha^\epsilon_h,T).
\]
Each of these deficits starts at zero and is nondecreasing along
the path.

At a pre-outgoing step of the corresponding sign let
\(b\) be the strip size and \(\alpha,\alpha'\) the old and
new lists for this sign. Monotonicity and interlacing give
\[
 0\le\sum_h\bigl(\min(\alpha'_h,T)-\min(\alpha_h,T)\bigr)
           \le \min(b,T).
\]
The \(b\) bound sums the coordinate increments. For the
\(T\) bound use \(\alpha'_{h+1}\le\alpha_h\) and telescope. At a step of the opposite sign the
second sum can only decrease. During the outgoing frame both
lists of parts decrease entry by entry by the tail interlacing,
and the reference lists stay fixed.

Write
\[
 \Xi_x=\mathcal L\beta_x/(pk).
\]
We bound the weight of
paths violating any of the following conditions at a late top
\(x\): a covector decreases positive mass by at most \(\Xi_x\),
a vector decreases negative mass by at most \(\Xi_x\), and the
last plane entry before either step is at least \(2B_u/3\).
Here \(u\) is the earliest of the preceding \(r\) vectors, which
is before the outgoing frame. These conditions imply the good
path inequalities. At a covector all plane entries are
nonnegative, and the complement's positive mass cannot
increase, so the plane loss is at most the positive mass loss.

For a complement weight \(\kappa\), denote the positive and negative masses by \(m^+=\sum_i(\kappa_i)_+\) and \(m^-=\sum_i(-\kappa_i)_+\). At a vector step the change in the sum of the entries is 
\[
 \Delta\!\sum_i\kappa_i=\Delta m^+-\Delta m^-,
\]
 the positive mass gain plus the negative mass loss. Interlacing bounds its
positive mass gain by the last plane entry before the step.
Reading that entry backwards through the last \(r-1\) vector
steps bounds it by the first entry just after \(u\). Geometric growth of the vector degrees gives
\[
 \text{first entry just after }u\ \le\!\sum_{x\le u,\ x\text{ a vector}} B_x\le 2B_u.
\]
Using the bound
\[
 2B_u\le4a/P_{j,u},
\]
the complement
gain is at most
\[
       \text{negative mass loss}+2B_u
                      \le \Xi_x+4a/P_{j,u}.
\]

\proofstep{The loss from the mass test.}
For a violation of the mass condition let \(U\) be the sum
of preceding degrees of the opposite sign. If there is any
mass of that sign, at least one such pre-outgoing degree has
appeared. Its largest degree is at least \(k\) and at most
\(4\beta_x/p\). The geometric sum gives
\[
 U\le5\beta_x/p.
\]

The total mass of the sign before this step is at most \(U\)
by the deficit inequality at \(T\ge S\). Hence for \(T\ge U\)
the deficit of that sign grows by more than \(\Xi_x\). This
includes negative mass losses at outgoing steps, since removing
an entry equal to \(S\) preserves the negative mass.

\proofstep{The loss from the last-entry test.}
For a violation of the last-plane-entry condition, put \(B=B_u\).
If \(t\) outgoing steps have passed at this cut, \(u\) is the
oldest of the
\[
 h=r-t
\]
most recent pre-outgoing vectors.
The sum of positive rows after the first \(h\) in the tail is
at most the sum of all vector degrees strictly before \(u\).
At each of those \(h\) vector additions, interlacing bounds the
positive sum below row \(i+1\) after the addition by the sum
below row \(i\) just before. Covectors and the later outgoing
trims can only decrease these sums. Thus the positive sum after
row \(h\) is at most \(B/(\rho-1)\). Row \(h\) has positive part
less than \(2B/3\).

For
\[
 B\le T\le13B,
\]
the \(h\) latest pre-outgoing
vectors contribute at least \((h-1)T+B\) to the reference sum.
In the clipped tail, the first \(h-1\) rows contribute at
most \((h-1)T\), row \(h\) at most \(2B/3\), and the rest at
most \(B/(\rho-1)\). Subtracting gives
\[
 \delta_+(T)\ge \big((h-1)T+B\big)
                 -\big((h-1)T+2B/3+B/(\rho-1)\big)\ge B/4
\]
by the stipulated growth of the vector degrees.

\proofstep{Converting exceptional events into lost weight.}
To convert these deficits into a weight loss, set
\(J_k(T)=\sum_{l=1}^{k-1}(T+l)^{-2}\). For \(b\ge0\),
\[
 \log D_k(b)=\int_0^\infty\min(b,T)J_k(T)\,dT,
\]
where the left side for real \(b\) denotes the logarithm of the
binomial product. The equality follows by integrating each
\((T+l)^{-2}\) separately.

At the terminal complement weight
\(\lambda\) in dimension \(K\), the positive and negative lists
bound its dimension by
\[
 \dim V_\lambda\le\prod_h D_k(\alpha_h^+)D_k(\alpha_h^-).
\]
Indeed each Weyl factor is bounded by
\[
 1+\frac{\lambda_i-\lambda_j}{j-i}
 \le\left(1+\frac{\max(\lambda_i,0)}{j-i}\right)
     \left(1+\frac{\max(-\lambda_j,0)}{j-i}\right).
\]
Multiplying gives the bound with \(D_K\), which is at most
\(D_k\) for each nonnegative argument. Therefore
\(\log(R/\dim V_\lambda)\) is at least the integral of
\((\delta_++\delta_-)J_k\) at the last cut. If the mass bound
fails, this integral is at least
\[
 \Xi_x\sum_{l=1}^{k-1}\frac1{U+l}
 \ge \frac{\mathcal L(k-1)}{9k}\ge\frac{\mathcal L}{16},
\]
using \(k\le4\beta_x/p\). If the plane cutoff fails, it is at
least
\[
 \frac B4\sum_{l=1}^{k-1}
     \left(\frac1{B+l}-\frac1{13B+l}\right)
 \ge\frac{11(k-1)}{104}\ge\frac{k}{16}.
\]
Here \(B\ge k\), and both estimates are independent of the
remaining path and its strip degrees.

At each cut the two lists of nonzero parts have lengths at
most \(m\), since each pre-outgoing strip introduces at most
one new part of its sign. The outgoing interlacing keeps that
bound after removing the heads. Each part is an integer in
\([0,S]\). The lists padded to length \(m\), and the time of
the cut, determine the whole shape. There are at most
\((1+S)^{2m^2}\) shape paths because the zero start is fixed.
Thus the exceptional weight is at most
\[
 R\exp\bigl(2m^2\log(1+2aN_s)-\min(k,\mathcal L)/16\bigr)
 \le R a^{-N_s^2}.
\]
The last inequality uses \(1+2aN_s\le a^2\), supplied by
\eqref{eq:plane-count-size}, and the bounds on \(k,\mathcal L\).
This proves \eqref{eq:plane-good-path-count} for each choice
of intervening digits with the same constants.
\end{proof}


\paragraph{Sum over outgoing rows and digits.}
The last gap before the first outgoing position exceeds \(L_1\).
Thus each pre-outgoing \(x\) has
\(P_{j,x}\ge p^L p^{L_1+1}/H_j\). The positions are distinct, and
\(B_x\le 2a/P_{j,x}\). Summing the geometric series and using
\(a/p^L<8pYk^2N_s\) gives
\[
 S\le \frac{4aH_j}{p^L p^{L_1+1}},\qquad
 4YS\le H_j .
\]

For each outgoing row of a fixed path, the congruence
\(d_i+\sum_h v_{i,h}\equiv n_i\pmod{H_j}\) has exactly
\(H_j^{r-1}\) solutions with \(0\le v_{i,h}<H_j\). Each coordinate
is uniform among these solutions since \(r\ge2\). The union
bound removes at most \(2rS H_j^{r-2}\) solutions on requiring
\(v_{i,h}<H_j-2S\) for all \(h\). Each surviving solution has
\(d_i+\sum_h v_{i,h}<rH_j\), and so supplies a permitted outgoing
degree. The count per row is at least
\(H_j^{r-1}(1-r/(2Y))\) when this factor is nonnegative.

\begin{proposition}[Rank after summing the coarsened families]\label{prop:coarse-plane-rank}
With the admissible frame degrees and path count
\eqref{eq:plane-good-path-count}, the coarsened generators
span a residue space \(\overline\Lambda\) with
\[
 \dim\overline\Lambda\ge q^{-rK}(1-Nr^2 k/Y)\prod_xD_k(b_x).
\]

\end{proposition}

\begin{proof}
\proofstep{Summation using the degree at the outgoing row.}
For fixed digits the number of decorated words depends only on the outgoing degree
of its own segment. The incoming padding changes with adjacent
outgoing degrees, but the independence and membership just
proved hold at every such degree. Hence the sum of products
of word counts over the outgoing choices factors across segments.
For the digit sum at a top between frames, dividing each of the
\(k\) exponents of a monomial of degree \(a'_x\) by \(P_{j,x}\)
gives the identity \(\sum_{\ell_x} D_k(n_x)=D_k(a'_x)\).

At a frame top there are exactly \(P_{j,x}^{k-1}\) digit choices
by the sum congruence. Its incoming path factor in the next
segment is \(D_k(W_{j+1,x}-m_0)\), so its full contribution is
at least
\[
 \frac{1-r/(2Y)}{H_j^K}\,
 (P_{j,x}H_j)^{k-1}D_k(W_{j+1,x}-m_0).
\]
This is where the complement dimension \(K\) in the exponent
arises. The product of the factors \(H_j^{-K}\) over all \(r\)
rows in each frame is \(q^{-rK}\).

For the other factors, \(E'_x+E_x\le5b_x/(kY)\).
The binomial product formula and Bernoulli's inequality give
\(D_k(a'_x)\ge(1-5/Y)D_k(b_x)\). For an incoming radix
\(P'=P_{j,x}H_j\), put \(B'=W_{j+1,x}-m_0\). The floor and
padding lose at most \(4kN_s\) in the coarse degree.

Since
\(a'_x/P'\ge 4Yk^2N_s\),
\[
 P'B'\ge a'_x-4kN_sP'\ge(1-1/(kY))a'_x,\qquad
 (P')^{k-1}D_k(B')\ge(1-1/Y)D_k(a'_x).
\]
For the latter comparison use \(P'(B'+i)\ge(1-1/(kY))(a'_x+i)\)
for each \(1\le i\le k-1\). If \(1-Nr^2k/Y\le0\), the
claimed lower bound is automatic. Otherwise every loss factor
is in \([0,1]\). Since \(a\ge32Y\), \(J\le N\), and \(rJ\le N\),
the product of all loss factors is at least
\[
 1-N(5+1+r/2+1/32)/Y\ge 1-Nr^2k/Y.
\]

The next bound counts independent residue symbols in the
coarse radices. The following paragraph then checks that the
same families are admitted in the original radices, as needed
for Lemma~\ref{lem:lift-plane-symbol}.

To see that the dimension counts add, recover the choices from
the exponent at each top. For a coordinate \(h\) in a shared
frame row, write its true exponent as
\[
 \ell_{x,h}+P_{j,x}\bigl(a_{x,h}+H_j b_{x,h}\bigr).
\]
The digit \(\ell_{x,h}\) is the remainder modulo \(P_{j,x}\).
The outgoing exponent \(a_{x,h}\) is less than \(H_j\)
by the coordinate condition, so division of the quotient by
\(H_j\) recovers \(a_{x,h}\) and \(b_{x,h}\). In a complement
coordinate \(b_{x,h}=0\). The sum of the outgoing exponents
recovers \(e_x\). At unshared tops there is no second quotient.
Hence multiplication induces an injective map on the monomial
tuples, and distinct choices of digit and outgoing degree give
disjoint monomial supports. By
\eqref{eq:plane-good-path-count} and the count just proved,
\begin{equation}\label{eq:plane-coarse-rank}
 \dim \overline\Lambda\ \ge\ q^{-rK}(1-Nr^2k/Y)\prod_xD_k(b_x).
\end{equation}

\end{proof}

\paragraph{Return to the original radices.}
\label{par:plane-rank-transfer}

The count so far used coarsened digit classes. Returning to the original radices determines the rank available for the actual symbols.

\begin{proposition}[Rank of symbols at the original radices]\label{prop:plane-original-symbol-rank}
In the setting of Proposition~\ref{prop:coarse-plane-rank},
uncoarsening and multiplication by the reserved determinants
give integral original symbols of degrees \(b_x-E'_x\)
satisfying the required prefix tests. Their span contains a direct summand of rank \(r_{\rm orig}\) satisfying the bound in \eqref{eq:plane-coarse-rank},
\[
 r_{\rm orig}\ge q^{-rK}(1-Nr^2k/Y)\prod_x D_k(b_x).
\]

\end{proposition}

\begin{proof}
It suffices to return the products of the counted words to
the original symbol module. Use \(P^0_{j,x}=p^{z_j-x}\),
\(F_j=P_{j,x}/P^0_{j,x}=p^{R_j}\). Euclidean division in the unique late segment gives
\[
 \ell_{x,h}=\ell^0_{x,h}+P^0_{j,x}u_{j,x,h},\qquad
 0\le\ell^0_{x,h}<P^0_{j,x},\quad 0\le u_{j,x,h}<F_j .
\]

To
uncoarsen a counted segment product, raise its plane and
complement factors to \(F_j\). Multiply each factor by the
corresponding coordinate monomial of exponents \(u_{j,x,h}\)
in the late rows. Multiply the plane factor also by the
ordinary determinants of that segment with powers \(g_{j,x}\).
All coefficients are integers. Indeed, for an integral factor \(A\) with integer coefficients, Frobenius gives
\[
 \overline{A(X)^{F_j}}=\overline{A}(X^{F_j})
\]
in the residue field, with coordinatewise powers on the right. Applying the original radices in every
segment and multiplying by the remaining digits \(\ell^0\)
therefore gives \(\overline{\mathcal D}\) times the coarsened
product. In the integral polynomial the total degrees are
the coarsened true degrees plus the exact determinant degrees
\(E_x\), hence equal to \(b_x-E'_x\).

\proofstep{The inequalities for the restored products.}
We check the original inequalities on these integral products.
The plane and complement letter supports are unchanged, since
the digits are placed only in late rows and the determinants
use earlier plane vectors. There are at most \(kN_s\) coordinate digits in a prefix and each has size at most \(F_j-1\). Their combined adverse change in separation is therefore at most
\[
 \sum_{x,h}u_{j,x,h}\le kN_s(F_j-1).
\]
An odd plane coordinate can lower the plane
minimum by at most its exponent, and an even complement
coordinate can raise the complement maximum by at most its
exponent.

The other coordinate types do not decrease the
separation. The determinant associated to a given cut lies
wholly before that cut, so it contributes \(g_{j,x}\) to the
minimum. All other determinants have nonnegative minimum
since they contain only vectors. Adding the determinant gain to the worst digit loss gives, at this cut,
\[
 \text{separation}\ge F_jm_0-kN_s(F_j-1)+g_{j,x}>m_*+g_{j,x}.
\]
For the strict inequality use the separation margin in \(m_0\) and
\(F_j>a/(8pYk^2N_sQ)>m_*\) from \eqref{eq:plane-count-size} and
the definition of \(L\).

\proofstep{The rank lower bound at a plane.}
For the absolute plane lower bound put \(M=b_u/P_{j,u}\) at
the earliest preceding vector \(u\). The count assumptions
include \(k\ge C_kN_s^2\log a\) with \(C_k\ge1088\). In
particular \(a'_u\ge9b_u/10\) by the bound on \(E_u+E'_u\).
Since \(u\) is pre-outgoing, the degree \(B_u\) used by the
count satisfies \(B_u\ge 9M/10-4kN_s\) for either an
incoming or an intervening row. Also \(M\ge8Yk^2N_s\ge
40kN_s\). The plane minimum before the cut in the uncoarsened
word is at least
\[
 F_j(2B_u/3-kN_s)\ge F_j M/2=\frac{b_u}{2P^0_{j,u}}.
\]

For the normal cutoff at \(x\), the extra digit degree is at
most \(kF_j\). The gap \(x-u>L_1\) and \(F_jP^0_{j,u}\le A\)
give \(kF_j\le a/(p^{L_1}P^0_{j,x})\). The term at an even
\(x\) from the coarsened normal bound is
\[
 4a/P^0_{j,u}\le a/(p^{L_1}P^0_{j,x}).
\]
Thus the total
above \(\mathcal L a/(pkP^0_{j,x})\) is bounded by
\(a/(p^{L_1}P^0_{j,x})\) at an odd top and by twice that at
an even top. These imply both of the original order bounds.

The original symbols produced this way have the thresholds,
normal cutoffs and degrees needed by the precision induction.
After the original radix substitutions, the total normal
degree at each top is bounded by
\(\mathcal L a/(pk)+10a p^{-L_1}+kQ\). Indeed there is one
late occurrence, the incoming contribution is tangential,
and the final digit costs at most \(kQ\). The same upper bound
holds on supports of sums, so the total over all tops is at
most
\[
 N(\mathcal L a/(pk)+10a p^{-L_1}+kQ)
\]
on the whole
retained lattice.

The determinant product has nonzero reduction. For each of
its factors the preceding \(r\) even positions are distinct
modulo \(n'\). An interval of length \(n'\) starting immediately
after the earliest of those positions contains all \(r\) even
positions of any given frame modulo \(n'\), each once. If the
last position were at least that far from the earliest, there
would be at least \(r\) later preceding positions in the same
segment. Therefore the determinant monomials remain distinct
under the positive-power substitutions and have coefficients
\(\pm1\).

Since the residue polynomial ring is a domain, multiplication by the nonzero determinant product has zero kernel,
\[
 \overline{\mathcal D}\sum_i a_iW_i=0\quad\Longrightarrow\quad\sum_i a_iW_i=0.
\]
Hence the
products counted in \eqref{eq:plane-coarse-rank} remain independent
after multiplying by \(\overline{\mathcal D}\). Complete their
independent reductions to a basis of the ambient fixed-degree
residue space. Lifting the additional basis vectors to the
integral monomial lattice gives a basis by a change of basis
with unit determinant. This gives an integral direct summand
of the indicated rank generated by original symbols.
\end{proof}

The preceding rank and lifting statements concern the symbol
with the reduced weights. In Proposition~\ref{prop:residue-plane-neighbor-test}
the source must have the full weight. The replacement below
retains the rank of the model family and supplies the integral
contractions used on the target side.

\paragraph{Restoring the bounded weights.}
\label{par:plane-middle-weights}
\hypertarget{source-plane-middle-weights}{}
At a top let \(\lambda=(\lambda_1,\ldots,\lambda_k)\) be its
true dominant weight and let \(b=\lambda_1-\lambda_2\). Remove
\((b,0,\ldots,0)\) to form the remaining dominant weight
\[
 \mu=(\lambda_2,\lambda_2,\lambda_3,\ldots,\lambda_k).
\]

Away from the tops take the removed weight to be zero and
\(\mu=\lambda\). At odd indices use the corresponding reversed
negative weights in the polynomial model. Write \(d_l(\gamma)\)
for the dimension of highest weight \(\gamma\) in dimension \(l\).
The local factor of the count for restoring weights is
\(d_{k-1}(\lambda_2,\ldots,\lambda_k)\) at a top and
\(d_k(\lambda)\) elsewhere. We construct integral polynomials
realizing exactly this count.

\begin{auxlemma}[Counting the tableaux at one index]\label{lem:bounded-tableau-count}
With the top weights \(\mu\) defined at the start of
paragraph~\ref{par:plane-middle-weights}, use tableaux of the
determinant-shifted shape. At an even top take tableaux with every entry in the first row
equal to zero. At an odd top take the reversed negative shape and exclude zero
from its alphabet. Writing \(n_s\) for the number admitted at position \(s\),
\[
 n_s=\begin{cases}d_{k-1}(\lambda_2,\ldots,\lambda_k),&s\text{ a top},\\d_k(\lambda),&s\text{ away from the tops}.\end{cases}
\]
Away from the tops every tableau is allowed. Distinct selected tableaux have distinct leading monomials over
every residue field.
\end{auxlemma}
\begin{proof}
Use a matrix of variables \(X=(X_{i,h})\), with rows and columns
numbered from zero. For a dominant weight \(\gamma\) the
polynomial realization after its determinant twist has shape
\(\gamma-\gamma_k(1,\ldots,1)\). Its integral basis consists of
products of initial-column minors indexed by semistandard
tableaux of this shape on the alphabet \(0,\ldots,k-1\).
For each tableau its leading monomial uses the diagonal of each
minor. Order the variables by column first and by increasing
row within a column, giving the earlier variable precedence.

In a minor the largest monomial uses the smallest available row
in the first column and then the smallest remaining row in each
successive column. For rows \(i_0<\cdots<i_{s-1}\), this gives
\[
 \operatorname{in}\det(X_{i_a,b})_{0\le a,b<s}
    =\prod_{b=0}^{s-1}X_{i_b,b}.
\]
The coefficient is one. If \(M_T\) is the diagonal monomial of the product for a tableau \(T\), then
\[
 \deg_{X_{i,b}}M_T=\#\{c:T(b,c)=i\}.
\]
Its weakly
increasing order reconstructs that row from these counts. These are the integral
standard monomial coordinates used for the component lattices.
\cite[Theorem II.2.16]{abw-schur}

Away from the tops take every tableau of \(\mu\) in the
corresponding signed realization. At an even top restrict to
those with first tableau row all zero. Removing that row
leaves the tableau of the smaller shape with alphabet
\(1,\ldots,k-1\), since every remaining box had a box above it.
This restriction counts
\[
 d_{k-1}(\lambda_2,\ldots,\lambda_k),
\]

by subtracting the common determinant shift. At an odd top
restrict to tableaux with no entry zero. The reversed negative
weight of \(\mu\) has two equal last entries. The determinant
shift leaves at least its last shape row empty. Dropping this
row and using the alphabet \(1,\ldots,k-1\) counts the dimension
of the reversed negative of \((\lambda_2,\ldots,\lambda_k)\),
which equals \(d_{k-1}(\lambda_2,\ldots,\lambda_k)\) by duality.
\end{proof}
Let \(n_2\) be the product of these counts and \(H_i\) the
resulting basis products in the independent block variables.

\begin{proposition}[Restoration of the bounded highest weights]\label{prop:restore-plane-model-rank}
Use the removed flat degrees and the \(n_2\) tableau products
\(H_i\) counted in Lemma~\ref{lem:bounded-tableau-count}. The products \(\theta(Q)H_i\) retain
independence modulo the uniformizer for every independent
family of flat polynomials \(Q\) of the prescribed degrees.
With the common factor \(C\) of \eqref{eq:plane-clean-factor}
and precision \(c'\ge\max(2(D_1+1),c_\star+o_S)\), the products
admit the stated clean approximations at the distinguished edge. With the bounded spread \(c_0\) and lower
degree \(a\), let \(f_{\rm true}\) and \(f_{\rm flat}\) be the retained fractions in the true and flat dimensions, respectively. In the indicated size range,
\[
 f_{\rm true}\ge(1-kc_0/a)^N f_{\rm flat}.
\]

\end{proposition}

\begin{proof}
\proofstep{Independence after restoring the removed weight.}
To prove independence after restoring the removed weight, let
\(\theta\) insert the first column at an even index and the
signed \((k-1)\)-row minors of the first \(k-1\) columns at an
odd index. This is the vector and covector polynomial model of
Section~\ref{sec:edge-transfer}. At an even top with removed
degree \(b\), the leading monomial of a degree-\(b\) monomial
in the flat variables inserts exponents \(\alpha_c\) at
\(X_{c,0}\). The exponent matrix of a chosen tableau has no
\(X_{c,0}\) for \(c>0\).

At each \(c>0\) the exponent of \(X_{c,0}\) in the leading product
is therefore precisely \(\alpha_c\). Recover the missing exponent from
the homogeneous degree:
\[
 \alpha_0=b-\sum_{c>0}\alpha_c.
\] At an odd top the leading
monomial of the cofactor omitting row \(c\) uses \(X_{t,t}\)
for \(0\le t<c\) and \(X_{t+1,t}\) for \(c\le t<k-1\).
In the selected tableau the entry in row \(t\) is at least
\(t+1\), so its exponent at \(X_{t,t}\) is zero. The diagonal exponents of the product read
\(d_t=\sum_{c>t}\alpha_c\) for \(0\le t<k-1\).

If \(b\) is the fixed flat degree, invert this triangular
reading by
\[
     \alpha_0=b-d_0,\qquad
     \alpha_c=d_{c-1}-d_c\ (1\le c\le k-2),\qquad
     \alpha_{k-1}=d_{k-2}.
\]
Thus the diagonal exponents and the degree recover the
full flat monomial. After recovering the flat exponents, subtract their contribution
to recover the leading exponent matrix of the tableau. At an
off-top position there is no contribution to subtract. Thus the
pairs of a flat monomial and a selected tableau product have
distinct leading monomials. Each has leading coefficient a unit,
including in the residue field. In a proposed nonzero relation
\[
 \sum_i\theta(Q_i)H_i=0,
\]
choose the largest of those leading
monomials among the nonzero coefficients. Its coefficient cannot
cancel with that of any other pair. This proves the assertion
over the residue field for arbitrary \(Q_i\) of the fixed degrees.

\proofstep{The three normalized edge lattices.}
Use the normalized edge coordinates on all three weights
\(\lambda,\mu,\lambda-\mu\). The normalization consists of
pullback to the component frames, the last-entry determinant
powers in the signed realizations, and the common edge powers
\(w^{A_r}\) for the respective weights. In each frame the common powers on the two factors have product
\[
 w^{A_r(\mu)}w^{A_r(\lambda-\mu)}=w^{A_r(\lambda)},
\]
by additivity of the edge exponent in the weight.

On a flat polynomial this change to edge
coordinates is the letter substitution followed by \(\theta\).
Indeed a vector transforms by the frame matrix \(M\), and the
cofactors by \(\det(M)M^{-t}\), which follows by multiplying
the adjugate identity by \(M\). The last-entry twist on the
covector cancels the determinant. The flat edge exponent is
one copy of the degree for each even block and zero for each
odd block. Write
\[
 D=\operatorname{diag}(w^{-1}I_r,I_{k-r})
\]
for the diagonal
part. After the common determinant cancels, the two letter maps are
\[
 wD=\operatorname{diag}(I_r,wI_{k-r})\quad\text{on vectors},\qquad
 D^{-t}=\operatorname{diag}(wI_r,I_{k-r})\quad\text{on covectors}.
\] The
other frame factors act by the same matrix identities, with
Frobenius at the corresponding index. Thus this normalization
pulls the edge coefficient lattices back to their integral
standard monomial models, including their reductions.

In the distinguished frame the model vectors \(S_b\) have
normalized polynomials \(B_b C\) with one common factor \(C\).
This is \(\operatorname{subst}(\mathrm{SPtrue})\) in the symbol
identity, with residue the product of \(F_*\) and the separate
determinant reserve on the digits. Each determinant uses distinct
rows and has a permutation monomial with unit coefficient, so
\(\overline C\ne0\). To obtain the clean source, expand \(B_b\)
by its segment presentation. Its nonnegative separation margin
allows each homogeneous pair to be expanded in the paired paths
of Section~\ref{par:plane-column-orders}. The common factor
\begin{equation}\label{eq:plane-clean-factor}
                 C=F_*\prod_s T_s^{kQ}
\end{equation}
supplies the precision reserve and the factors for the digits. Write \(S_{b,o}\) for the standard-edge pullback of the
normal degree-\(o\) component of \(C B_b\). Apply the path
and digit constructions to this component.

In the flat source
lattice \(R_e\) on the chosen edge, the resulting vector
\(F_{b,o}\) has source order at most \(o\) and satisfies
\[
                 F_{b,o}-w^o S_{b,o}\in w^{o+c'}R_e .
\]
The insertions commute with
the projection to the full true multidegree because \(C\)
has the reserved degrees. Divide by \(w^o\) in the displayed bound,
then sum over the finitely many \(o\) bounded in
Section~\ref{par:plane-rank}. Thus
\[
  F_b:=\sum_o w^{-o}F_{b,o}\in S_b+w^{c'}R_e .
\]
Here \(S_b=\sum_o S_{b,o}\). The competing-edge conditions
of the construction give \(F_b\in w^{c'}R_{e'}\) on every
other edge out of the same source. Restore the weights by multiplying
with the inverse edge transports of the \(H_i\). They are integral
on the chosen edge, so the same congruence holds in the restored
lattice there.

If a combination with coefficients \(c_{b,i}\in O\) is
zero modulo \(w\) in the normalized edge lattice, then in the
residue polynomial model
\[
 0=\sum_i\theta\left(\overline C\sum_b\overline{c_{b,i}}\overline{B_b}\right)
                    \overline{H_i}.
\]
All polynomials inside \(\theta\) have the full common flat
multidegree. By block independence, for each \(i\) the coefficient polynomial
vanishes. Cancellation takes place in the residue polynomial ring,
a domain, and then independence of the \(\overline{B_b}\) gives
\[
 \overline C\sum_b\overline{c_{b,i}}\overline{B_b}=0
 \ \Longrightarrow\ \sum_b\overline{c_{b,i}}\overline{B_b}=0
 \ \Longrightarrow\ c_{b,i}\in wO\quad\text{for every }b .
\] In particular primitive combinations remain primitive
in the actual source lattice.

\proofstep{Pairing the restored model family.}
For the pairing one can therefore test the restored model family.
An integral target test evaluates the difference between the clean
and model combinations in \(w^{c'}O\). If \(c'\ge 2(D_1+1)\), this
is contained in \(p^{D_1+1}O\), since \(w^2=p\). A model evaluation
of valuation at most \(D_1\) has the same valuation on the clean
combination. On competing source edges restoration costs at most
\(o_S\) powers of \(w\), by the source vertex bound for the remaining
weights. Thus \(c'\ge c_\star+o_S\) gives the source conductor there.

At each top the spread \(\lambda_2-\lambda_j\) for \(j\ge2\)
is at most \(c_0\), with \(b\ge a\) in the plane count.
The Weyl factors involving the first entry give
\[
 d_k(\lambda)\le d_{k-1}(\lambda_2,\ldots,\lambda_k)D_k(b)
                  (1+c_0/a)^{k-1}.
\]
This follows term by term from
\(\lambda_1-\lambda_j+j-1\le (b+j-1)(1+c_0/a)\).
The retained fraction of the true dimension is therefore at
least the flat fraction multiplied by \((1-kc_0/a)^N\), using
\((1+t)^{-(k-1)}\ge1-(k-1)t\) for \(t\ge0\). The product of
the off-top dimensions is included exactly.
\end{proof}

The pairing of two lifted symbols is computed from the unit parts of factorials in the prescribed digit classes. The identities below give the
precision for that pairing in the proof of
Proposition~\ref{prop:residue-plane-neighbor-test}.
\paragraph{Factorials and divided derivatives in a digit class.}
\label{par:plane-digits}
\hypertarget{source-plane-digits}{}
We record two uniform identities used in constructing the signed-chain
pairing. Write \(P=p^g\) and \(0\le\rho<P\), and use the residue
class \(\rho+P\beta\) for integers \(\beta\ge0\). Congruences in
this paragraph are coefficientwise over \(\mathbb Z_{(p)}\) and
therefore hold over the coefficient DVR as well.

For \(M\ge0\), let \(U(M)=\prod_{1\le v\le M,\,p\nmid v}v\),
with \(U(0)=1\). Removing the multiples of \(p\) from a factorial
and repeating \(g\) times gives
\begin{align*}
 (\rho+P\beta)!&=\beta!\,p^{A_\rho+\beta B_g}
       \prod_{h=0}^{g-1} U(b_h+P_h\beta),\\
 b_h&=\left\lfloor\frac\rho{p^h}\right\rfloor,\quad
 P_h=p^{g-h},\quad A_\rho=\sum_{h=0}^{g-1}\left\lfloor\frac{b_h}{p}\right\rfloor,
 \quad B_g=\frac{P-1}{p-1}.
\end{align*}
For each \(h\) put \(C_h=U(P_h)\). Since \(b_h<P_h\), splitting
into full blocks followed by the last partial block gives
\begin{align*}
 U(b_h+P_h\beta)={}&U(b_h)C_h^\beta\,
 \prod_{j=0}^{\beta-1}\prod_{\substack{1\le v\le P_h\\p\nmid v}}
                (1+P_hj/v)\,
 \prod_{\substack{1\le v\le b_h\\p\nmid v}}(1+P_h\beta/v).
\end{align*}
The two products on the right together form a unit. In a group
of coordinates with the same \(g\), the powers of \(p\) and \(C_h\)
contribute a constant depending on the total \(\sum_a\beta_a\)
and the fixed residues. The remaining unit is the product of
the displayed products of \(1\) plus a multiple of \(p\).
This gives the factorial factorization with the exact degree
of the quotient monomials as the argument of the constant.

\begin{lemma}[Factorial units and divided derivatives at fixed digits]\label{lem:plane-digit-units-and-cuts}
On a fixed digit class \(\rho+P\beta\), with \(P=p^g\),
\(g\ge0\), \(0\le\rho<P\), \(\beta\ge0\), and \(c\ge1\), either
sign of the variable unit factor in the factorial is congruent
modulo \(p^c\) to a product-binomial polynomial of total degree
at most \(2(c-1)\). A system of divided derivatives of
orders \(M_a\) admits the quotient expansion
\eqref{eq:plane-digit-derivative-expansion} with the following bound when the common radix is \(P\),
\[
 \sum_a\kappa_a\le\frac1P\sum_a M_a+c-1.
\]
 In a signed cut a divided derivative of order
\(|M|\) loses at most \(|M|\) on a vector letter and zero
on a covector letter. The factorial formula and expansions
include the radix-one case.
\end{lemma}

\begin{proof}
\proofstep{Uniform polynomial descriptions of units and inverse units.}
For \(c\ge1\), the remaining unit and its inverse both have binomial
polynomial approximations to precision \(p^c\) of total
degree at most \(2(c-1)\). To prove it, replace the leading
factor \(p\) in each \(P_hj/v\) and \(P_h\beta/v\) by a variable
\(z\). For either sign, expand the factors to degree \(c-1\)
as polynomials in \(z\) modulo \(z^c\).

In a full block at \(j\)
the coefficient of \(z^d\) is a scalar in \(\mathbb Z_{(p)}\)
times \(j^d\). Multiplying over \(j\) expresses the coefficient
of \(z^d\) as a combination of sums
\[
 \sum_{0\le j_1<\cdots<j_l<\beta}j_1^{d_1}\cdots j_l^{d_l},
 \qquad d_i\ge1,\quad \sum_i d_i=d.
\]
Induct on the number \(l\) of indices. For the induction step,
fix the last index \(j_l\). The sum over its predecessors is
an integer-valued polynomial in \(j_l\) of degree at most
\((d-d_l)+(l-1)\). Multiplication by \(j_l^{d_l}\) gives
degree at most \(d+l-1\). Summing from zero to \(\beta-1\)
raises the degree by at most one, since
\[
 \sum_{j<\beta}\binom jt=\binom{\beta}{t+1}.
\]

This proves the degree bound \(d+l\) and integrality of
the values at nonnegative integers. A partial block contributes an ordinary
polynomial of degree \(d\) in \(\beta\) at \(z^d\). The coefficient
at total order \(d\) over all blocks and coordinates is thus a
polynomial of total degree at most \(2d\), taking values in
\(\mathbb Z_{(p)}\) at nonnegative integer tuples. Write \(P(\boldsymbol\beta)\) for this coefficient polynomial and
\(\Delta_i\) for forward difference in coordinate \(i\). Its Newton
expansion is
\[
 P(\boldsymbol\beta)=\sum_{|\mathbf h|\le2d}
    (\Delta_1^{h_1}\cdots\Delta_m^{h_m}P)(0)
    \prod_{i=1}^{m}\binom{\beta_i}{h_i},
\]
where \(m\) is the number of coordinates. Each coefficient is an
integral linear combination of values of \(P\), hence lies in
\(\mathbb Z_{(p)}\). Now
put \(z=p\). This proves the approximation uniformly in the
number of coordinates and quotient exponents, also when the
fixed residues and powers \(g\) differ by group. If \(g=0\)
its product is empty and the same formulas give unit one.

\proofstep{Divided derivatives and their factorial exponents.}
For the divided derivative, write \(\partial^{[M]}X^n=\binom nM
X^{n-M}\), with the derivative zero for \(n<M\). Use
\[
 (1+T)^P=1+T^P+pE(T),\qquad \deg E\le P-1,
 \qquad E\in\mathbb Z[T].
\]
Then \(\binom{\rho+P\beta}{M}\) is the coefficient of \(T^M\) in
\((1+T)^\rho(1+T^P+pE(T))^\beta\). At order \(\ell<c\) in \(pE\),
the part with \(T^{Pj}\) from the other factor has coefficient
\[
 p^\ell[T^{M-Pj}](1+T)^\rho E(T)^\ell\,
       \binom{\ell+j}{\ell}\binom\beta{\ell+j}.
\]
This uses the binomial identity for selecting \(\ell\) marked
factors and then \(j\) others. The support of this expression
requires \(j\le M/P\) and \(M-Pj\le\rho+\ell(P-1)\).

Put \(t=\lceil(M-\rho)/P\rceil\) and \(\rho'=\rho-M+Pt\), so
\(0\le\rho'<P\). With \(\kappa=\ell+j\), the support condition
gives \(\kappa\ge t\). Therefore the corresponding divided
derivative term has the form
\begin{equation}\label{eq:plane-digit-derivative-expansion}
 \begin{aligned}
 &c_{\ell,j}X^{\rho'}
       \left(X^{\kappa-t}\partial^{[\kappa]}X^\beta\right)(X^P),\\
 &c_{\ell,j}=p^\ell[T^{M-Pj}](1+T)^\rho E(T)^\ell
                                      \binom{\ell+j}{\ell}.
 \end{aligned}
\end{equation}
All constants are integral and independent of \(\beta\).
Terms with \(\beta<\kappa\) vanish. For multiple coordinates
multiply these expansions and keep \(\sum_a\ell_a<c\), so
\[
 \sum_a\kappa_a\le\sum_a M_a/P+c-1.
\]
The discarded coefficients
are divisible by \(p^c\), also after grouping equal monomials.
This is the required derivative expansion inside the old residue
class, with the new residues \(\rho'_a\) fixed by \(M_a\).
The case \(P=1\) has \(E=0\) and gives the ordinary divided
monomial derivative exactly.

\proofstep{Effect of a derivative on a signed cut.}
In the same argument a prefix lower bound is the least Laurent
exponent for the conjugates of \(\operatorname{diag}(T,1,\ldots,1)\)
acting on the letters in the cut. For any rank-one projection
\(P_1\), write this matrix as
\[
 A_T=1+(T-1)P_1.
\]

The inverse is \(1+(T^{-1}-1)P_1\). A vector coordinate in the
cut has least order at least zero and a covector coordinate
has least order at least \(-1\). A divided derivative of
order \(|M|\) in one letter loses at most \(|M|\) in the vector
case and zero in the covector case. To see this, write the
substitution of that letter as \(x\mapsto B_Tx\). Expand the
polynomial after substitution at \(x+B_T^{-1}h\) and take the
coefficient of \(h^M\). Thus the Laurent orders of the entries satisfy
\[
 \operatorname{ord}_T\!\big((B_T^{-1})_{ij}\big)\ge\begin{cases}-1,&B_T^{-1}=A_T^{-1}\quad(\text{vector}),\\0,&B_T^{-1}=A_T^t\quad(\text{covector}).\end{cases}
\]
For a letter
outside the cut the Jacobian is the identity. The coefficient
estimate holds for each conjugate and hence for the cut bound.

Adjoining vector letters to a cut composes its substitution
with polynomial substitutions on those separate letters.
It can only increase the least order. Removing them composes
with inverse substitutions and loses at most the sum of their
degrees by homogeneity in each removed letter.  They give the local cutoff losses
when the binomial approximation is implemented by divided
derivatives and multiplication by monomials of the same degree.
\end{proof}

To obtain the covector in Proposition~\ref{prop:residue-plane-neighbor-test},
we compare the actual pairing with its model on prescribed cut
subspaces. The corrected pairing estimate below uses the same original multidegree as the source lift in
Lemma~\ref{lem:lift-plane-symbol}.
\paragraph{The Fock metric and the signed cuts.}
\label{par:plane-gram}
\hypertarget{source-plane-gram}{}
We give the pairing argument in its signed polynomial model.
Normalize the valuation by \(v(p)=1\). For a monomial with
exponent \(\alpha\), write \(\alpha!=\prod_x\alpha_x!\) and
\(h_\alpha=v(\alpha!)/2\). Use the bilinear form and its
associated coefficient bound
\[
  [F,G]=\sum_\alpha F_\alpha G_\alpha\,\alpha!,
  \qquad \nu(F)=\min_\alpha \bigl(v(F_\alpha)+h_\alpha\bigr).
\]
Then \(v([F,G])\ge \nu(F)+\nu(G)\). In true multidegree
\(d=(d_i)\), set
\[
 n_0=\frac12 v\!\left(\prod_i d_i!\right),\qquad
 K_d=\sum_i r\bigl(1+\log_p(d_i+r)\bigr),\qquad
 e_d=r^2\bigl(1+\log_p(2\textstyle\sum_i d_i+r)\bigr).
\]

In this multidegree \(2n_0-K_d\le v(\alpha!)\le2n_0\).
Indeed the difference is the valuation of a product of
multinomial coefficients. Each multinomial is a product of
at most \(r\) binomial coefficients with upper entries at most
the letter degree. A binomial with upper entry \(n\ge1\) has
valuation at most \(1+\log_p n\), by the base-\(p\) carry
formula, which follows from Legendre's formula for factorials.
The same argument for a split \(\alpha=A+M\) gives
\[
   v(A!)+v(M!)\ge 2n_0-2K_d,
   \qquad \text{when \(A+M\) has multidegree \(d\)}.
\]

\begin{lemma}[Projectors for signed cuts in the factorial norm]\label{lem:plane-signed-projectors}
For every prescribed cut and threshold in fixed multidegree
\(d\), there is a Fock-self-adjoint rational projector onto
the signed cutoff subspace, whose loss in \(\nu\) is at most
\(e_d\). For nested or disjoint cuts write \(P_i\) for their projectors and \(P\) for the stipulated product. Then
\[
 P^2=P,\qquad \operatorname{im}P=\bigcap_i\operatorname{im}P_i,
\]
with loss at most \(e_d\) per cut.
It preserves multidegree and the monomial exponents outside
the union of those cuts.
\end{lemma}

\begin{proof}
\proofstep{The projector for a signed cut.}
Using the signed, generically conjugated Laurent order, the space attached to a cut \(S\) is
\[
 U(S,b)=\{F:\operatorname{ord}_S(F)\ge b\}.
\]
In fixed
multidegree it is the sum of the \(\mathrm{GL}_r\)-types
for the action on \(S\) whose least weight \(\lambda_r\) is
\(\ge b\). Indeed all torus exponents in a type are at least
\(\lambda_r\) in the selected coordinate, and the conjugates
of that torus test detect the lowest weight in every nonzero
copy. More explicitly the intersection of the conjugate
weight-cut spaces is invariant, and the copies it contains
are exactly the stated types. The type sum has a Fock
self-adjoint projector \(\Pi\), defined over \(\mathbb Q\).

By the trace formula and divided-power estimate for type
projections proved above, its matrix on the coefficient lattice satisfies
\[
 v_p(\Pi_{\beta\alpha})\ge-e_d\qquad(\text{all }\alpha,\beta).
\]
In applying the
estimate the coordinate span on the cut is at most \(\sum_i d_i\),
and it suffices to use \(2\sum_i d_i\) in the Hom weight bound.
Zero degree and \(r=1\) give scalar projections on the affected
monomial spaces. The trace formula works for the sum of types
with the same bound by the ultrametric inequality.

For the matrix entry \(\Pi_{\beta\alpha}\), self-adjointness
gives
\[
 \beta!\Pi_{\beta\alpha}=\alpha!\Pi_{\alpha\beta}.
\]

Consequently the two numbers
\(v(\Pi_{\beta\alpha})+h_\beta-h_\alpha\) and
\(v(\Pi_{\alpha\beta})+h_\alpha-h_\beta\) are equal.
Their common value is the mean of the two unweighted
entry valuations, each bounded below by \(-e_d\).
For the coefficient of \(x^\beta\) in \(\Pi F\) the
ultrametric inequality therefore gives
\[
 v((\Pi F)_\beta)+h_\beta
    \ge\min_\alpha\{v(F_\alpha)+h_\alpha-e_d\}.
\]
Taking the minimum in \(\beta\) yields the claimed norm bound.

For nested cuts the smaller-cut projector commutes with
the action on that cut and with the action on its complement.
It therefore preserves every isotypic component for the
larger cut and commutes with its type projector. For disjoint
cuts the projectors act on separate factors. A product over
such cuts is self-adjoint, fixes their intersection and loses
at most \(e_d\) per cut. It preserves multidegree and each
exponent outside the union of the cuts. This proves the projection and gives the product bound used below.

\end{proof}

\begin{paramdefinition}[The true and model lattices, cuts and signs]\label{hyp:plane-gram-data}
The block letters \(B_j\) in each frame \(j\) have specified
signs and total block degree \(b_j\). The shared letters \(X_t\),
indexed by \(t=(j,m)\), \(0\le m<r\), are vectors. Let
\(P_t=p^{g_t}\) be their radices. The model has the same
block letters and three vector letters \((L_t,W_t,T_t)\)
over each shared letter. Its collapse \(\phi\) sends them
respectively to \(X_t,X_t^{P_t},X_t\), coordinate by coordinate.
The exponents on \(L_t,W_t,T_t\) will be denoted by
\(\rho_t,\beta_t,M_t\). An admissible monomial has
\[
 \begin{gathered}
 0\le\rho_{t,a}<P_t, \\ M_{j,0}=0, \\ |M_{j,m}|\le b_j,
\end{gathered}
\]
and exact collapsed multidegree \(d\).

\smallskip
The cuts are defined as follows.
A true link cut contains \(B_j\) and the shared letters
\(X_{j,m'}\) for \(m'>m\). Denote it by \(A_{jm}\);
we use \(m+2\le r\). The corresponding model link cut
\(U_{jm}\) contains \(B_j\), the letters \(T_{j,m'}\)
for \(m'>m\), and all the \(W_{j-1,m'}\), where the frame
indices are cyclic.

A model chain cut \(S_{ji}\) contains a
nested subset of \(B_j\) and all the \(W_{j-1,m'}\).
Thus the link cuts alone, the chain cuts alone, and the
union of the chain cuts with the model link cuts are each
families of nested or disjoint cuts. Set
\[
 D_l=J(r-1)e_d,\qquad D_m=e_d\sum_j n_j,
\]
where \(J\ge1\) is the number of frames and \(n_j\) the
number of chain cuts at \(j\). The model total degree is
at most \(\sum_i d_i\) on admissible monomials because
each \(P_t\ge1\).

Use a diagonal unit \(\gamma(\alpha)\) on the true
monomials. Its two signs admit integral binomial polynomial
approximations of total degree at most \(6c\) modulo \(p^c\).

Suppose that \(\nu(F)=0\) in true multidegree \(d\), that
the true link orders of \(F\) are at least \(\ell_{jm}\),
and that each digit part of \(F\) has chain orders at least
\(\theta_{ji}\). The digit part for a fixed residue keeps the
block exponents and reads on \(W_t\) the quotient exponent
\[
 m_t=\rho_t+P_t\beta_t\quad\longmapsto\quad \beta_t=(m_t-\rho_t)/P_t,
\]
with no \(L_t,T_t\) exponents.
The pairing under consideration is \([F,\gamma\,\phi(Z)]\),
where \(\gamma\) acts coefficientwise. Use the inequalities
\begin{equation}\label{eq:plane-gram-reserves}
 \begin{split}
    c&>4(D_m+D_l+3K_d),\\
    s_{ji}&\ge r\, b_{j-1}/P_{j-1,m}+13rc
                     \qquad(0\le m<r).
 \end{split}
\end{equation}
\end{paramdefinition}
The family of test polynomials is indexed by primitive vectors of the true lattice. The two signs of the diagonal unit are treated by Lemma~\ref{lem:plane-digit-units-and-cuts}. The reserves are inequalities on the integer conductor chosen for the small-prime calculation; its specialization is made in Section~\ref{sec:host-edge-tests}.


\begin{proposition}[An integral polynomial with prescribed orders at the model cuts]\label{prop:plane-corrected-gram-polynomial}
Use the lattices, orders and sign approximations in
Definition~\ref{hyp:plane-gram-data}, including \(\nu(F)=0\)
and the reserves \eqref{eq:plane-gram-reserves}.
Then there is an integral model polynomial \(Z_*\) of collapsed
degree \(d\), with model link orders at least
\(\ell_{jm}-30c\) and chain orders at least
\(\theta_{ji}-s_{ji}-12c\), such that
\[
 v([F,\gamma\phi Z_*])-n_0+K_d
                               \le D_l+D_m+4K_d.
\]
The order and valuation bounds are uniform over the coefficients,
digit classes and radices with the data of Definition~\ref{hyp:plane-gram-data}.
\end{proposition}

\begin{proof}
\proofstep{Projecting a monomial to the link cuts.}
Take a monomial \(x^\alpha\) attaining \(\nu(F)=0\) and
project it to the true link cuts at levels \(\ell_{jm}\).
The result \(Y\) has \(\nu(Y)\ge h_\alpha-D_l\), and
\[
 [F,Y]=F_\alpha\,\alpha!
\]
has valuation \(h_\alpha\).
Apply an integral binomial approximation \(g\) of \(\gamma^{-1}\)
of degree at most \(6c\) to its coefficients. A factor
\(\binom{\alpha_x}{a}\) acts as
\(x^a\partial_x^{[a]}\). Its loss in a signed cut is at
most \(a\), by the cut estimates for the derivative and
multiplication. This bounds the total link loss by \(6c\).
The change in the pairing from \([F,Y]\) to
\([F,\gamma(gY)]\) has valuation at least
\[
 c+h_\alpha-D_l>h_\alpha.
\]

Ordinary coefficients of \(gY\) have valuation at least
\(h_\alpha-D_l-n_0\). Multiplication by
\[
 p^{\lceil n_0+D_l-h_\alpha\rceil}
\]
therefore makes \(gY\)
an integral polynomial \(G\) in the original variables. It has link orders at least
\(\ell_{jm}-6c\) and its twisted pairing has valuation at
most \(n_0+D_l+1\).

\proofstep{Lifting \(G\) in the prescribed plane.}
To lift \(G\), first put each \(X_t\) entirely in \(L_t\).
Within frame \(j\), order its block letters first and then
the shared letters in decreasing order of \(m\). Each
\(A_{jm}\) is a prefix of this order. Use the integral basis
of signed path words proved in Section~\ref{sec:column-lifts},
now testing the least torus weight at the prefixes.
In a positive column the least exponent selects a
subset avoiding the torus coordinate whenever \(h<r\).
In a nonpositive column it selects a dual subset
containing the coordinate whenever \(h<r\). Thus if
\(h_b\) are the prefix heights,
\[
 \lambda_r=
 \sum_{b\ {\rm positive}}\mathbf1_{h_b=r}
       -\sum_{b\ {\rm nonpositive}}\mathbf1_{h_b<r}.
\]
The Pieri decomposition at the successive prefixes
selects exactly the paths whose displayed sums obey all
the lower bounds. The associated words span the
intersection of the cuts. Since they are a
subset of the integral basis, the expansion of an integral
polynomial in this intersection uses integral coefficients.

In a path word, a column meeting at least one block
letter moves all its shared occurrences with \(m>0\) from
\(L_{j,m}\) to \(T_{j,m}\). Keep \(L_{j,0}\) and every
occurrence in a column disjoint from the blocks. Each
column uses a given shared letter at most once. Count the columns by one of their block occurrences. Each chosen occurrence lies in only one column, so
\[
 \#\{\text{moved columns}\}\le\#\{\text{block occurrences}\}\le b_j.
\]
This gives the tangent degree bounds.
On a moved column the new cut \(B_j\cup\{T_{j,m'}:m'>m\}\)
acts on exactly the old cut occurrences. An unmoved
column has zero order there. Its contribution to the
prescribed lower bound was at most zero.

In fact for a
positive column the height at the cut is bounded by the
number of its occupied vector letters before the cut,
which is \(<r\) since the column meets only shared letters.
For a prefix, a moved column retains its signed cut
order and an unmoved column contributes zero, at least
its previous order. Minimum Laurent orders add on products, giving the required cut
bound for the whole word. For the coefficient valuation the relevant
identity is
\[
      \nu_{\rm coeff}(AB)=\nu_{\rm coeff}(A)+\nu_{\rm coeff}(B)
                 \qquad(A,B\ne0).
\]
Indeed, divide each factor by its minimum uniformizer power. Their
reductions are nonzero polynomials over the residue field, so their
product has nonzero reduction. The injective renaming on each
column preserves that column's \(\nu_{\rm coeff}\), and
hence preserves the valuation of the product word.
Collapsing \(T_{j,m}\) back to the corresponding shared
letter restores each column separately.

Perform the operation coefficientwise outside each frame.
Since the renamed letter sets are disjoint, that operation
commutes with the cut substitutions in the other frames.
For a remaining shared variable, divide its exponent by its radix.
The lift and collapse act on the monomial by
\[
 a_t=P_tq_t+r_t,\quad 0\le r_t<P_t,\qquad
 L_t^{r_t}W_t^{q_t}\:\longmapsto\: X_t^{r_t+P_tq_t}=X_t^{a_t}.
\]
The tangent exponents stay fixed throughout. To see that it preserves
the link bounds, group by the \(L\)-exponents. These letters
are outside the new link cuts, so coefficient extraction
preserves the bounds. The added monomial in the \(W\) letters
has nonnegative order at every model link cut, since those
letters are vectors. The integral model polynomial \(Z\)
therefore has admissible support, \(\phi(Z)=G\), and
\(\operatorname{ord}_{U_{jm}}Z\ge\ell_{jm}-6c\).

\proofstep{Correction to impose the chain cuts.}
For an admissible model exponent \(\zeta\), let \(A\) be
its collapsed exponent without \(T\), and let \(M\) be its
tangent exponent on the true shared letters. Thus
\(\phi(x^\zeta)=x^{A+M}\). The digit factorial identity
gives
\[
       A! M! = \kappa_B\,u_\zeta\,\zeta!,
       \qquad
       \eta_\zeta=\gamma(A)u_\zeta ,
\]
where \(u_\zeta\) is the product of the digit units.
The constant \(\kappa_B\) depends on the individual label and
tangent exponents and on the model multidegrees, which we fix
in each block of the next projection. Explicitly it is the
product of the digit constants for each \(W_t\)-degree divided
by the product of the label factorials. The letter \(T_t\)
contributes the same factorial on both sides.

\proofstep{Diagonal correction factors.}
Choose diagonal multipliers \(g_+\) and \(g_-\) approximating
\(\eta_\zeta\) and its inverse modulo \(p^c\), with
integral values and cut cost \(12c\) on cuts excluding \(L\).
For fixed label exponents, substitute
\(\rho_{t,a}+P_t\beta_{t,a}\) into the binomial approximations
to \(\gamma^{\pm1}\). Apply forward differences in \(\beta\) to the integer-valued polynomial \(A(\beta)=\binom{\rho+P\beta}{k}\). Its expansion is
\[
 A(\beta)=\sum_{j=0}^k(\Delta^j A)(0)\binom{\beta}{j},\qquad (\Delta^j A)(0)\in\mathbb Z.
\]

Multiplying this approximation by the degree-\(6c\) digit
unit approximation yields the cost \(12c\). Extracting
coefficients of \(L\) leaves every such cut intact, so the
estimate holds after adding the label classes.

In each model multidegree let \(\Pi_m\) be the product
projector on chain orders at least \(\theta_{ji}-s_{ji}\).
Use its linear extension over all multidegrees. It fixes
the exponents on \(L,T\), preserves the admissible support
conditions, and commutes with all model link projections.
Each multidegree component of a polynomial satisfying a
link bound still satisfies it, by the homogeneous cut action.
 Set
\[
 \Omega=g_-\Pi_m g_+.
\]

Its chain orders are \(\ge\theta_{ji}-s_{ji}-12c\).
Any model link bound is lowered by at most \(24c\).
For \(\zeta,\zeta'\) in the same fixed block of \(\Pi_m\),
\[
 \lvert v(\zeta!)-v(\zeta'!)\rvert\le K_d.
\]

Only the block and \(W\) exponents can vary. Their degrees
are bounded by the corresponding true degrees, so the
multinomial estimate applies without a second factor for
the fixed \(L,T\) exponents. The factorial-weighted projection
bound proves that \(p^{N_2}\Omega\) preserves integrality
on admissible support, with
\(N_2=\lceil D_m+K_d/2\rceil\).

\proofstep{Precision of the corrected pairing.}
For a fixed tangent exponent \(M\) put
\[
         F_M=\gamma^{-1}\,\partial^{[M]}(\gamma F).
\]
Every coefficient of \(F_M\) has valuation \(\ge -n_0\):
the unit multipliers and the integral divided derivative
preserve the original coefficient bound. At a fixed label
\(\rho\), its digit part has a decomposition \(Q+E\) such
that
\[
 v(E_\beta)\ge c-n_0
\]
and
\(\operatorname{ord}_{S_{ji}}Q\ge\theta_{ji}-s_{ji}\)
for all \(j,i\). To obtain it, start with the digit parts
of \(F\). In each of the two twist operations replace the
unit multiplier by its degree-\(6c\) approximation. Each
replacement error has coefficient valuation at least
\(c-n_0\), and each approximation loses at most \(6c\) in
the chain orders, using the affine binomial expansion.
Between the two multipliers use the divided derivative
expansion of the preceding subsection on each shared letter.

The divided-derivative expansion takes the input label to the residue of \(\rho-M\)
coordinatewise. The error coefficients have valuation at
least \(c-n_0\). The approximate derivatives and all
multiplying monomials have integral coefficients, so errors
already made retain this valuation. Only derivatives on
the \(W\) letters in the cut can lower its order. Their
total loss is bounded by
\[
    \sum_{m=0}^{r-1}
       \left(\frac{|M_{j-1,m}|}{P_{j-1,m}}+c\right)
    \ \le\ s_{ji}-12c
\]
by \eqref{eq:plane-gram-reserves} and \(r\ge1\).
This produces \(Q,E\) with the asserted relation.

\proofstep{Comparing factorial terms in the corrected pairing.}
For an admissible monomial the factorial identity and the
coefficient formula for the divided derivative give
\[
 [F,\gamma\,\phi(x^\zeta)]
       = \kappa_B\,\eta_\zeta\, (F_M)_A\,\zeta! .
\]
In the model multidegree of \(\zeta\), multiply \(Q\) by
the fixed label and tangent monomial, then keep that
multidegree. The resulting polynomial \(Q_B\) still has
the required chain orders because \(L,T\) are outside
those cuts. On that block the difference
\[
    g_-(\zeta')\,\eta_{\zeta'}(F_M)_{A'}-(Q_B)_{\zeta'}
\]
has valuation at least \(c-n_0\). The pairing with
\(Q_B\) is unchanged by \(\Pi_m\).

To estimate the remaining
error from a coefficient of \(\Pi_m x^\zeta\), use its bound
\(h_\zeta-h_{\zeta'}-D_m\). From the split factorial
bound applied to \(\zeta\) and \(\zeta'\), and the equality
of their \(\kappa_B\), we have
\[
 v(\kappa_B)+h_\zeta+h_{\zeta'}\ \ge\ 2n_0-2K_d .
\]
The error terms in comparing the evaluation before and
after \(\Pi_m\) thus have valuation at least
\(c+n_0-2K_d-D_m\), including the original monomial term.
Applying \(g_+\) at the input preserves this error estimate.
The factor \(g_+(\zeta)g_-(\zeta)-1\) has valuation at least
\(c\), and the original evaluation has valuation at least
\(n_0-K_d/2\) by \(\nu(F)=0\) and the true factorial bound.
Linearity and integral coefficients of \(Z\) now give
\begin{equation}\label{eq:plane-gram-stability}
 v\bigl([F,\gamma\phi(\Omega Z)]-[F,\gamma\phi(Z)]\bigr)
               \ \ge\ c+n_0-2K_d-D_m .
\end{equation}

\proofstep{The available precision after the conductor loss.}
The strict conductor inequality in \eqref{eq:plane-gram-reserves}
makes \eqref{eq:plane-gram-stability} larger than
\(n_0+D_l+1\). Hence the integral test
\(Z_*=p^{N_2}\Omega Z\) has valuation of its pairing at
most \(n_0+D_l+1+N_2\). Its model link and chain orders are
\(\ge\ell_{jm}-30c\) and \(\ge\theta_{ji}-s_{ji}-12c\),
respectively. Its collapse has true multidegree \(d\) and
integral coefficients, hence \(\nu(\phi Z_*)\ge n_0-K_d\).
In particular
\[
 v([F,\gamma\phi Z_*])-n_0+K_d
                 \le D_l+D_m+4K_d .
\]
Here \(K_d\ge1\) absorbs the two roundings. The estimate is uniform and needs no error polynomial.
The chain and link losses and the valuation gain are
uniform over all input coefficients, label classes, and
radices at the specified degrees and cuts.
\end{proof}

The Gram test must apply to the actual span of allowed
operators. We verify the prefix bounds on that span before
using the chosen-plane projection in
Proposition~\ref{prop:residue-plane-neighbor-test}.
\paragraph{Prefix orders on the exact operator span.}
\label{par:plane-orbit}
\hypertarget{source-plane-orbit}{}
In the distinguished edge coordinates the flat generators of the model
are exactly \(B_b\,\mathrm{SPR}\). Indeed pulling back each row of the
true frame determinant by the edge substitution gives its designated
plane row, so the product \(\mathrm{SPtrue}\) pulls back to
\(\mathrm{SPR}\). Every monomial in this common factor uses only
even plane coordinates. We work over the fraction field in their
linear span and apply words in \(X_{sa}\partial_{sb}\), using each
top \(s\) at most \(c_0\) times. Conjugation by an invertible linear
substitution within each letter takes a row operator to a linear
combination of row operators at the same letter. In matrix-unit notation the substitution by \(g\) acts on a row operator through
\[
 gE_{ab}g^{-1}=\sum_{c,d}g_{ca}(g^{-1})_{bd}E_{cd}.
\]
Each term acts at the same top, so the word allowance persists in edge coordinates.

List the tops after the incoming frame of segment \(j\) in increasing
order as \(x_0,\ldots,x_{h+r-1}\), where \(h\) is the number strictly
between the frames. The last \(r\) positions are the outgoing frame,
with shared index
\[
 m=r-1-(u-h)
\]
at position \(x_u\).
Write \(B_u\) for the true degree of the oldest of the \(r\) even
predecessors of \(x_u\), divided by its radix in this segment. Put
\[
 s_u=\lfloor B_u/(Yk^2)\rfloor.
\]

For a joint complement type,
let \(k_u\) denote its largest weight at the prefix of the first
\(u\) of the listed tops. In this paragraph frame subscripts on
\(B_u,s_u,k_u\) are implicit. Define \(T_u\) as the largest of
the following applicable quantities:
\[
 \frac{B_u}{2},\quad k_{u+1}+m_*+s_u\quad(x_u\text{ even}),\quad
 k_{u-1}+m_*+s_u\quad(u>0,\ x_{u-1}\text{ odd}).
\]
The chain and link thresholds for the plane Gram argument are
\begin{equation}\label{eq:plane-chain-link-thresholds}
   T_u-s_u/8\quad(0\le u\le h),\qquad
   T_u-3s_u/8\quad(h<u<h+r),
\end{equation}
respectively. Its separate chain
slack is \(3s_u/8\). We show that these thresholds work for
every choice of low digits, quotient multidegrees, joint
complement type and coefficient of that type.

\begin{proposition}[Prefix thresholds on the full permitted operator span]\label{prop:plane-operator-prefix-orders}
On the fraction-field span of the flat generators \(B_b\,\mathrm{SPR}\)
and the allowed words in \(X_{sa}\partial_{sb}\) with at most
\(c_0\) operators at each top, the chain and link thresholds from \eqref{eq:plane-chain-link-thresholds} are
\[
 \begin{cases}T_u-s_u/8,&0\le u\le h,\\T_u-3s_u/8,&h<u<h+r.\end{cases}
\]
 They hold after taking any indicated digit class,
quotient multidegree and joint complement type, with the same
allowance \(s_u\) for the selected coefficient.
\end{proposition}

\begin{proof}
\proofstep{Expansion into segment generators.}
First expand a generator in its presentation by products of
segment pairs \(A_jG_j\) in the plane and complement
coordinates. Do this before selecting a coefficient. The external
label \(l\) at each top is strictly below the segment radix there.
Take homogeneous components at the prescribed true degrees of
\(B_b\) in this expansion. Renamed products of homogeneous segment
factors have definite true degrees, and the sum of the components
in other degrees is zero. Thus the degree bound below applies on
each retained product. Expand \(\mathrm{SPR}\) into monomials
\(X^\gamma\).
In a coordinate of radix \(P\) and prescribed residue \(\rho\), the term can survive extraction only if
\[
 l+\gamma\equiv\rho\pmod P,\qquad (l+\gamma-\rho)/P\in\mathbb Z_{\ge0}.
\]
Apply this test in every coordinate.
In the matching class the quotient after multidegree projection
is a tensor product. Its complement factor is the projected
product of the \(G_j\) with their positions renamed.

The
complement is absent on incoming rows and the complement label
has zero quotient. Its plane factor, before multidegree projection,
is
\[
        X^b\,\phi\!\left(\prod_j \chi_j A_j\right).
\]
Here \(b\) is the plane quotient of \(l+\gamma\) by the low
radices, coordinate by coordinate. It vanishes on odd letters.
The variable map \(\chi_j\) sends intervening, outgoing, and
incoming letters respectively to the block letters,
\(L_{j,m}\), and \(W_{j-1,m}\). If \(P_{\rm low}\) is the low
radix on the preceding frame, the exponent multiplier at an incoming
row is
\[
           P_{\rm coarse}=P_{\rm low}\,P_{j-1,m},
                 \qquad P_{j-1,m}=p^{z_j-z_{j-1}}.
\]
This proves the displayed formula after collapse, since \(\phi\)
raises \(W_{j-1,m}\) to the \(P_{j-1,m}\)-th power.
These maps use the original frame radices after the transport in
Section~\ref{par:plane-rank}.

If the complement projection vanishes, the extracted plane
coefficient is zero. Otherwise its type at every prefix has
largest weight at most the upper Laurent order of the corresponding
prefix of \(G_j\). The other complement factors are outside the
cut, and the multidegree projection commutes with its action.
An isotypic projection with larger first weight would annihilate
the vector by the conjugate torus criterion. The original segment
inequalities now give
\[
 \operatorname{ord}_{<x_u} A_j\ge T_u.
\]

For an even letter use the complement through that letter. For a
preceding odd letter use its complement before it and the plane
through it. On segment support this last plane prefix is precisely
the prefix before \(x_u\). Its oldest even predecessor and hence
the available margin are unchanged in traversing the odd letter.
The absolute plane lower bound gives the first term of the maximum.
The extracted coefficient is a scalar multiple of the projected
plane factor. All these conditions describe linear spaces.

\proofstep{Transport of the prefix cuts by relabeling.}
For the chain cut \(S_{ju}\), the renaming \(\chi_j\) carries the
prefix before \(x_u\) to all incoming \(W\)'s and the first \(u\)
intervening letters with their original signs. Other segments
have no variables in this cut. Thus
\[
 \Pi=\prod_j\chi_j A_j
\]

has order at least \(T_u\). Restrict to a desired true quotient
multidegree of \(X^b\phi\Pi\) before applying the shared digit
extraction. On \(\Pi\) this restriction selects the model
multidegrees for which the weighted sum of letter degrees after
collapse has the prescribed value. Their projections commute with
the chain action. Expand the resulting polynomial by monomials
\(L^M\) of the label letters, which are outside the cut. Each
coefficient, a polynomial in block and \(W\) letters, retains the
order.

For each shared coordinate \(t\), divide the exponent in
\(X^b\phi(L^M)\) by \(P_t\):
\[
     b_t+M_t=\rho_t+P_t a_t,\qquad 0\le\rho_t<P_t,\quad a_t\ge0 .
\]
Thus the coefficient contributes only to the residue \(\rho_t\)
and acquires the fixed carry \(W_t^{a_t}\). For a tested residue \(\rho'_t\), the carry factor after extraction is therefore
\[
 \mathbf1_{\rho'_t=\rho_t}W_t^{a_t}.
\]
Apart from these carries the extra monomial
is the block part of \(X^b\). All affected coordinates of this monomial are vectors.
They have nonnegative least weight and multiplication preserves
the lower bound. This proves the required chain order for every
shared residue.

A true link at \(x_u\) with \(u>h\) sees the intervening letters
and the earlier outgoing rows. The original prefix also sees the
incoming rows. In a homogeneous component of \(A_j\), let their
total degree be \(d_W\). Let \(D(T)\) be the rank-one one-parameter torus
used for the original Laurent substitution.

To omit the
incoming variables \(W\) from this cut, substitute \(D(T)^{-1}W\)
in the transformed polynomial. A monomial of degree \(d_W\) accumulates \(d_W\) inverse weights, each either zero or \(-1\). Its Laurent order in this substitution is therefore
\[
 -\sum_{a=1}^{d_W}\varepsilon_a\ge-d_W,\qquad \varepsilon_a\in\{0,1\}.
\]
Therefore
\[
 \operatorname{ord}_{\rm prefix\ without\ incoming}(A_j)
 \ge \operatorname{ord}_{\rm original\ prefix}(A_j)-d_W .
\]
Coefficient extraction in the now untransformed \(W\)
preserves this bound for each remaining polynomial. Under collapse those incoming monomials go to frame
\(j-1\), outside this true link. Here \(J\ge2\) follows from the
frame span bound \(G'\) and \(3G'\le n'\), since one frame would
recur only after \(n'\) positions. All factors from other
segments and from \(b\) that meet this link use vectors there.
Projection to the required multidegree preserves the lower bound.

Only an incoming degree \(d_m^W\) with
\[
 P_{j-1,m}d_m^W\le (d_{\rm pl}(j-1,m)-|b_{j-1,m}|)_+
\]

can contribute. At the row \(y_m\) of that preceding frame put
\(t_m=z_{j-1}-y_m\), \(g=z_j-z_{j-1}\). Removing the quotient
of the external monomial leaves at most the original generator
degree \(2a\) at true scale. Indeed after multiplying the
remaining quotient by \(p^{t_m}\) the degree contributed by
the segment products is bounded by their prescribed true degree,
which is at most the flat degree. The \(t_m\) are distinct and
nonnegative. Consequently
\[
 d_W\le\sum_m 2a p^{-g-t_m}\le 3a p^{-g}.
\]

If \(y\) is the oldest even predecessor of the current outgoing
row, there are at least \(r\) intervening even tops before it is
tested. The first intervening top is more than \(L_1\) after
\(z_{j-1}\). Hence \(y\ge z_{j-1}+L_1+1\), and
\(a\le B_u p^{z_j-y}\). The host spacing gives
\(p^{L_1+1}\ge48Yk^2\); the size bound gives
\(B_u/(Yk^2)\ge1\), since \(B_u\ge a/Q\).
The permitted incoming degree is therefore bounded by
\[
 d_W\le \frac{3B_u}{p^{L_1+1}}
     \le \frac{B_u}{16Yk^2}\le s_u/8 .
\]
The last step uses \(B_u/(Yk^2)\le2s_u\).
Subtracting this degree from the full prefix order \(T_u\)
gives the initial link bound.

\proofstep{Admitting row operators in the order estimate.}
It remains to allow row operators. For one coordinate of low radix
\(P\), write its true power as \(v+Pq\), \(0\le v<P\).
Write \(x^v F(X)\) with \(X=x^P\). The two derivative
cases follow from the chain rule:
\[
 \partial_x\!\bigl(x^vF(x^P)\bigr)=
 \begin{cases}
   x^{v-1}(v+P X\partial_X)F(X),&v>0,\\
   x^{P-1}P\partial_X F(X),&v=0 .
 \end{cases}
\]
For multiplication the two cases are
\[
 x\,x^v F(X)=
 \begin{cases}
 x^{v+1}F(X),&v+1<P,\\
 X F(X),&v=P-1 .
 \end{cases}
\]
At \(P=1\) the second case applies to \(v=0\). The same
formulas at the second digit extraction act on the incoming
\(W\) quotient, or on the block coordinate at radix one. On an exponent \(v+Pq\), the Euler operator \(x\partial_x\) acts by \(v+Pq\). Thus on the quotient coordinate \(y\) in the residue class \(v\) its action is
\[
 v+P\,y\partial_y.
\]

Thus each multiplication or differentiation before extraction costs
at most one in a lower cut order that meets the letter. Under
conjugation the torus on a coordinate has weights \(0,1\) on a
vector and \(0,-1\) on a covector, with opposite signs on the
derivatives. The Euler operator also costs at most one, since one
of its two factors has nonnegative least weight. On a disjoint cut
these operations preserve the order.

For a complement operation, insert the joint decomposition of the
input before the operation and project the output. At each prefix \(I\), the upper-order estimate for an
operator at \(x\) allows a transition from \(k_I\) only to
\[
                          k'_I\le k_I+\mathbf1_{x\in I}.
\]
Complement coefficients commute with the plane cuts. Thus
this transition changes the maximum defining \(T_u\) by
at most one for a complement operation. It can
occur only for a top at or before \(x_u\) in that segment.
A plane operation can affect the shared \(W\) cut even if its
letter belongs to the previous frame, but that row occurs in the
incoming part of the segment, also before \(x_u\).
This accounts for both occurrences of a shared row, including
across the cyclic endpoint.

\proofstep{The exact operator span.}
If \(a_u\) operators in the word act at such tops up to the tested position \(u\), induction therefore bounds their cumulative loss by 
\[
 \underbrace{2+\cdots+2}_{a_u\,\text{operators}}=2a_u.
\]
 The
cut and type estimates for a single half-operation apply in all
input digit classes and degrees. The induction therefore survives
both quotient-degree changes and the sum over input types. For one cut, at most \(N_1c_0\) operators can touch its letters.
The host inequalities give
\[
 a\ge16N_1c_0Yk^2Q,
\]
and
\[
 B_u\ge a/Q
\]
implies
\[
 s_u\ge16N_1c_0.
\]

For a word let \(m_I\) be the number of operators meeting
the cut. The induction allows at most \(2m_I\) in the
chain and at most \(4m_I\) in the link, with both occurrences
of a shared row included. The budget \(m_I\le N_1c_0\) and
\(16N_1c_0\le s_u\) in particular give
\[
    \operatorname{ord}_{\rm chain}\ge T_u-s_u/8,\qquad
    \operatorname{ord}_{\rm link}\ge T_u-s_u/8-s_u/4 .
\]
The first loss is charged from the starting chain threshold
and the second from the initial link threshold.
These link thresholds are positive since \(T_u\ge B_u/2\)
and \(s_u\le B_u/(Yk^2)\). Linear combinations over the fraction field
preserve these inequalities. This proves the membership used for
the selected plane polynomial in Section~\ref{par:plane-flat-test}.
 \label{cor:plane-orbit-1}
\end{proof}

The preceding order and Gram results first give a test at
the flat weight. We use them to select a nonzero plane
component, realize its symbol, and restore the weight. This
supplies the detecting value used in
Proposition~\ref{prop:residue-plane-neighbor-test}.
\paragraph{Choosing the plane component.}
\label{par:plane-flat-test}
\hypertarget{source-plane-flat-test}{}
Let a flat source polynomial \(H\) in the full \(k\) coordinates
in the exact operator span of Section~\ref{par:plane-orbit} have
integral coefficients in the chosen edge frame after scaling.
Take \(H\) primitive there. At each letter use the radix
specified by its frame position. Fix the labels \(\rho\)
and the degrees of the plane and complement quotients using a
monomial of \(H\) with unit coefficient. In this class, write
\(H_{\alpha,\beta}\) for the coefficient at plane quotient
\(\alpha\) and complement quotient \(\beta\). Write
\((\alpha_0,\beta_0)\) for the two quotients of the selected
monomial. All the letter positions are distinct. Coordinatewise, radix \(P\), digit \(\rho\), and quotient \(q\) give
\[
 m=\rho+Pq,\qquad 0\le\rho<P.
\]

Taking \(q\) from the plane or complement quotient as appropriate therefore recovers every true exponent. Denote the assembly of a plane polynomial
\(A\) and a complement polynomial \(N\) by
\(\operatorname{Asm}_\rho(A,N)\): multiply the fixed label
monomial by the two polynomials with variables raised to
their letter radices. Their degrees are fixed by the chosen
class. The factorial identity of Section~\ref{par:plane-digits}
gives the exact formula
\begin{equation}\label{eq:plane-slice-factor}
 \begin{split}
 [H,\operatorname{Asm}_\rho(A,N)]
 &= b_\rho
     \sum_{\alpha,\beta} H_{\alpha,\beta}\,
       A_\alpha\,u_{\rm pl}(\alpha)\alpha!\,
       N_\beta\,u_{\rm no}(\beta)\beta!.
 \end{split}
\end{equation}
Here \(b_\rho\) depends on the labels and the quotient
multidegrees, and \(u_{\rm pl},u_{\rm no}\) are the products
of the factorial units in the two sets of coordinates.
The unit \(u_{\rm pl}\) is the twist \(\gamma\) in the
plane Gram estimate. Both signs have the required uniform
approximations. The formula is obtained monomial by
monomial, so it holds for the whole class.

\begin{lemma}[A nonzero plane polynomial and a complementary test]\label{lem:plane-split-polynomial}
Let \(H\) be an integral primitive flat polynomial in the
operator span of Proposition~\ref{prop:plane-operator-prefix-orders}.
Use the digits and quotient degrees of one unit-coefficient
monomial to determine the label class in \eqref{eq:plane-slice-factor}.
Then there
are a nonzero plane polynomial \(F\), an integral complementary
test \(N_0\), and a scalar \(k_0\ne0\) such that
\eqref{eq:plane-slice-value} and \eqref{eq:plane-slice-depth}
hold for every integral plane polynomial \(A\) at the selected
quotient degrees. In particular the scalar in \eqref{eq:plane-slice-depth} satisfies
\[
 v(k_0)\le V_d-n_0+D_n.
\]
The indices of the complementary types and their order losses are chosen before \(A\).
\end{lemma}

\begin{proof}
\proofstep{The projections onto the complementary factors.}
On the complement take the product projections onto one
isotypic type at each of its chain cuts. These cuts are
nested within a frame and disjoint across frames.
The products over all choices of types sum to the identity.
For \(Y=\sum_\beta H_{\alpha_0,\beta}x^\beta\), there is a
product \(E\) with \(\nu(EY)\le\nu(Y)\). Otherwise
\(\nu(\sum_E EY)>\nu(Y)\) by the ultrametric inequality,
contrary to the sum being \(Y\). In particular \(EY\ne0\).
For this \(E\), choose an exponent \(\tau\) attaining
\(\nu(EY)=q_1\), and put \(u=v(\tau!)/2\). The plane
polynomial \(P\) has coefficients
\[
 P_\alpha=\left(E\sum_\beta H_{\alpha,\beta}x^\beta\right)_\tau .
\]

The selected plane coefficient has \(v(P_{\alpha_0})=q_1-u\), and
\[
 q_1\le\nu(Y)\le v(\beta_0!)/2,\qquad
 q_3:=\nu(P)\le q_1-u+v(\alpha_0!)/2 .
\]
The second bound for \(q_1\) uses the unit coefficient.
Choose \(z\) with \(v(z)=q_3\), possible because the
coefficients have valuations in \(\frac12\mathbb Z\)
and \(w^2=p\). The normalized plane polynomial \(F=P/z\)
then has \(\nu(F)=0\). The orbit calculation gives the signed chain and link
orders for precisely this coefficient \(P\) of the projected
complement. Scalar multiplication by \(z^{-1}\) preserves them.

\proofstep{Scalar orders in the selected plane.}
Let
\[
 n_{0,{\rm no}}=\tfrac12v(\prod_i d_{{\rm no},i}!)
\]

for the complement degrees, and let \(D_{\rm no}\) be an
integer loss bound in the factorial norm for \(E\).
It is the sum of the cut projector losses, rounded up.
There is a scalar \(s\) of valuation
\[
 n_{0,{\rm no}}+D_{\rm no}-u.
\]

The complement polynomial \(s E x^\tau\) is integral
by the weighted projection bound. It has the upper torus
weight bound of the chosen type at each cut. Apply a binomial
approximation to \(u_{\rm no}^{-1}\) with integral coefficients
of degree \(\le 2(c_P-1)\) modulo \(p^{c_P}\).
This gives an integral complement test \(N_0\). In the upper
weight convention a diagonal binomial of degree \(a\) raises
the bound by at most \(a\), by applying the Laurent estimate
with the reciprocal torus variable. Thus the upper bounds of the chosen types increase by at most the permitted allowance \(4c_P\).

If one uses the exact inverse twist in place of the
approximation, Fock self-adjointness of \(E\) in
\eqref{eq:plane-slice-factor} gives exactly
\(k_0[F,\gamma A]\) with
\[
 k_0=b_\rho s\,\tau!z.
\]

Let \(V_d=v(\prod_i d_i!)\) be the full factorial valuation.
For each exponent in full multidegree let \(\kappa_{\rm fl}\)
bound its factorial deficiency from \(V_d\).
For example \(\sum_i k\lfloor\log_p d_i\rfloor\) works,
with a zero term if \(d_i=0\), by counting carries.
When \(A\) is integral the complement approximation gives
\begin{align}
 v\bigl([H,\operatorname{Asm}_\rho(A,N_0)]-k_0[F,\gamma A]\bigr)
       &\ge c_P+V_d-\kappa_{\rm fl}, \label{eq:plane-slice-value}\\
 v(k_0)&\le V_d-n_0+D_n, \label{eq:plane-slice-depth}
\end{align}
where \(n_0\) belongs to the selected plane degrees and
\[
 D_n=D_{\rm no}+
        n_0-\tfrac12v(\alpha_0!)+
        n_{0,{\rm no}}-\tfrac12v(\beta_0!).
\]
Indeed the factorial of the selected monomial has
valuation \(v(b_\rho)+v(\alpha_0!)+v(\beta_0!)\le V_d\).
Using \(v(z)=q_3\le v(\alpha_0!)/2+v(\beta_0!)/2-u\)
in the expression for \(v(k_0)\) gives
\eqref{eq:plane-slice-depth}. Assembly preserves coefficient
valuation bounds, giving \eqref{eq:plane-slice-value}
from the full factorial lower bound.

\end{proof}

\paragraph{Using a realised flat symbol.}
Apply the plane Gram estimate with conductor \(c_P\), which
is an integer here. Its chain and link losses \(12c_P\) and
\(30c_P\) also satisfy the looser allowances \(24c_P\) and
\(60c_P\) of the flat construction. In the signed-chain
hypothesis the slack \(26rc_P\) suffices for the \(13rc_P\)
in \eqref{eq:plane-gram-reserves}. Let \(A=\phi Z_*\)
and write \(D_b=D_m+D_l+3K_d\). The construction before
the last simplification in the Gram estimate gives
\[
 v([F,\gamma A])\le n_0+D_l+1+\lceil D_m+K_d/2\rceil
                 \le n_0+D_b .
\]
Here the integral conductor data have \(K_d\ge1\).

If \(c_P>D_n+D_b+\kappa_{\rm fl}\), the error in
\eqref{eq:plane-slice-value} is above the depth of the
main term. We therefore obtain
\[
 0\ \le\
 v\left(\frac{p^{\kappa_{\rm fl}}}
              {\prod_i d_i!}[H,\operatorname{Asm}_\rho(A,N_0)]\right)
 \ \le\ D_n+D_b+\kappa_{\rm fl}.
\]
The left inequality uses coefficient integrality and the
factorial deficiency termwise. No error polynomial is
needed for this application of the Gram estimate. The bounds and decomposition below apply uniformly
to all choices of the primitive source.

The detecting value must have a depth bound common to all
primitive source combinations. We choose this bound from the
weight data before choosing a detecting polynomial. Lemma~\ref{lem:lift-plane-symbol}
then lifts the resulting flat symbol to the integral polynomial lattice.
\paragraph{A depth fixed before the choice of a source.}
\label{par:plane-uniform}
\hypertarget{source-plane-return}{}
Let
\[
 L=Nk^2\log_p a.
\]
Here \(a\) denotes the lower bound for
the flat degrees on the tops, as in Section~\ref{par:plane-rank}.
By the inequality on the selected parameters \(a\ge\max(p,4N,k,2)\), and the full
degrees are at most \(2a\). There are at most \(N\) true letters
in the signed model: their positions in the cycle are distinct tops.
Each quotient degree is at most \(2a\) divided by its low radix.
These statements hold before the digits, projections or unit
monomial of the source have been chosen.

\begin{proposition}[An integral lift of the detecting symbol at uniform depth and order]\label{prop:plane-uniform-symbol-return}
Under the hypotheses on the frames and the parameter sizes, let \(X\) be the normal degree bound in the proof and put
\[
 D_X=\lceil49Nk^2\log_p a\rceil,\qquad t_F=\lfloor X\rfloor+Nc_0.
\]
 For each primitive polynomial
in the exact model span there is a detecting flat symbol
of order at most \(t_F\) at depth \(D_X\), realizable
at the selected plane with competitor conductor \(c_P\).
Both \(D_X\) and \(t_F\) are fixed by the weight and size
data before choosing the primitive polynomial.
\end{proposition}

\begin{proof}
\proofstep{The common bounds before choosing a source.}
In \eqref{eq:plane-slice-depth} take for \(D_{\rm no}\) the sum of
the integral projector exponents. At a complement cut of total
degree \(D\), the trace projector bound gives
\(\binom K2\lfloor\log_p(2D+K-1)\rfloor\), interpreted as zero
for a scalar action. The estimate also holds in the factorial
norm by self-adjointness. There are at most \(2N\) such cuts,
since the segment with \(h+r\) listed tops has \(h+r+1\) cuts.
Now
\[
 2D+K\le4Na+k\le a^3;
\]
hence \(D_{\rm no}\le6L\).
Counting carries at each true letter bounds the plane and
complement factorial deficiencies by \(3Nr\log_p a\) and
\(3NK\log_p a\), respectively. Their halves and the
projector loss thus give \(D_n\le9L\).
The full factorial loss \(\kappa_{\rm fl}\) is at most
\(Nk\log_p(2a)\le 2Nk\log_p a\le20L\).

Similarly \(e_d\le4k^2\log_p a\) and
\(K_d\le4Nk\log_p a\), since \(2\sum_i d_i+r\le a^3\).
There are at most \(N\) links and at most \(N\) chain cuts,
so \[
 D_b=D_m+D_l+3K_d\le20L.
\] We may take the integer
\[
                         D_X=\lceil49L\rceil\le50L .
\]
Choose the absolute constant \(C\) in \(D_p\ge CL\) before
taking the host sufficiently large. The conductor hypothesis gives
\[
 c_P\ge c_\star\ge8(D_p+1)>\max\{2D_X,4D_b\},
 \qquad D_X\le50L,\quad D_b\le20L .
\]

For the second reserve in \eqref{eq:plane-gram-reserves}, write
\(D=z_j-a_{j,1}\) and \(P=p^{z_{j+1}-z_j}\). Intervening letters
have distinct distances \(m\ge D+L_1+1\) from \(z_j\).
On a chain cut in segment \(j+1\) the oldest even predecessor
lies at or after \(a_{j,1}\). These two facts give
\[
 b_j\le 2a\sum_{m\ge D+L_1+1}p^{-m}
     \le 3a p^{-D-L_1-1},
 \qquad B_u\ge \frac{a}{p^D P}.
\]
Set \(x=B_u/(Yk^2)\), so \(s_u=\lfloor x\rfloor\).

At the tested letters \(B_u\ge a/Q\) and the host size bound
give \(s_u\ge100r^2 c_P\); in particular \(x\le2s_u\).
The spacing bound \(p^{L_1+1}\ge48rYk^2\) follows from the
host spacing with \(r\le Y\) and \(N\ge1\). Hence
\begin{equation}\label{eq:plane-cut-budget}
 \frac{rb_j}{P}\le \frac{3rB_u}{p^{L_1+1}}
       \le\frac{x}{16}\le\frac{s_u}{8},\qquad
 26r c_P\le\frac{s_u}{4},\qquad 128c_P\le s_u .
\end{equation}
The last two inequalities use \(r\ge2\). Adding the first two costs in \eqref{eq:plane-cut-budget} gives
\[
 rb_j/P+26r c_P\le3s_u/8
\]
on the same cut.

\proofstep{An integral symbol from the model polynomial.}
\label{par:plane-reassembly}
Use the integral model polynomial \(Z_*\) given by the Gram
estimate, and the complement test \(N_0\). Give each monomial
in the model the segment positions of its letters. A block letter
has its intervening position, \(L_{jm}\) and \(T_{jm}\) both have
the outgoing position of that shared row, and \(W_{jm}\) has
its incoming position in the next segment. Keep segment indices
on variables even when their positions in the cycle agree. Place
the complement variables only at the intervening and outgoing
positions. These renamings commute with the signed action on a
prefix and its inverse image. This follows on the variables;
the frame rows are all even and translation by \(n'\) preserves
parity. Coincidence of the label and tangent positions causes
no loss for this assertion.

\proofstep{Prefix orders in the model plane.}
For the renamed plane polynomial the prefix before \(x_u\)
has order at least
\[
                        t_u=T_u-s_u/2-60c_P .
\]
For \(u\le h\) the inverse image is the chain cut. The initial lower bound for the order is \(T_u-s_u/8\).
The two losses in the Gram construction are at most \(3s_u/8\) and
\(24c_P\), respectively. The resulting order is at least
\[
 (T_u-s_u/8)-(3s_u/8+24c_P)=T_u-s_u/2-24c_P\ge t_u .
\]

For \(u>h\) the inverse
image consists of the model link and the outgoing label letters
at earlier rows. These extra letters are vectors. Enlarging the
torus action to them multiplies the previous transform by a
commuting substitution with nonnegative powers, so the lower
bound
\[
 T_u-3s_u/8-60c_P
\]
suffices. For the prefix through an
intervening letter the chain cut at \(u+1\) gives \(t_{u+1}\).
The renamed complement polynomial has upper orders at most
\(k_u+4c_P\) and
\[
 k_{u+1}+4c_P
\]
before and through \(x_u\).

These bounds hold on the full polynomials over \(O\). To use them
segment by segment, take homogeneous components in each
segment letter on the plane side. The complement letter degrees
are already fixed. In one component all the permitted coefficients
in a particular segment, after fixing monomials in the other
segments, belong to the same integral submodule given by the
prefix inequalities. In fact taking such coefficients commutes
with the group of the specified segment. This submodule is
saturated in the finite free module of the given degrees:
its inequalities are vanishing conditions on Taylor or Laurent
coefficients over the fraction field. The quotient is torsion free,
hence free over the DVR.

For the fixed segment degrees, let \(M_j\) be the full
polynomial lattice and \(S_j\subset M_j\) the saturated
submodule cut out by its prefix bounds. Choose an integral
retraction
\[
 r_j:M_j\longrightarrow S_j,\qquad r_j|_{S_j}=\mathrm{id}_{S_j},
\]
possible since the quotient
is free. Fixing all monomials outside \(j\) lands in \(S_j\),
so the endomorphism \(r_j\) in factor \(j\) fixes the component
\(A\). These endomorphisms commute and hence
\[
                   A=(\bigotimes_j r_j)A
                           \in\bigotimes_j S_j.
\]
Applying the same map to a monomial expansion produces
a finite integral sum of products of permitted factors.
The identical argument on the complement introduces no
denominators there either.

\proofstep{Extracting factors with the true degrees.}
Write \(A_j,G_j\) for the factors of a resulting pair. Both have
the designated letter degrees. For an even intervening \(x_u\) take \(M=k_{u+1}+4c_P\),
the upper order of \(G_j\) through the letter. The plane
order before it is at least \(t_u\), and
\[
 t_u-M\ \ge\ m_*+s_u/2-64c_P\ \ge\ m_* .
\]

For an odd intervening \(x_u\) take \(M=k_u+4c_P\)
before the letter. The plane order through it is at least
\(t_{u+1}\) and
\(t_{u+1}-M\ge m_*+s_{u+1}/2-64c_P\ge m_*\).
The lower bounds follow from the respective inequalities
for \(T_u,T_{u+1}\).
The absolute lower bound before each letter is
\(B_u/2-s_u\), since \(T_u\ge B_u/2\).
These are exactly the two parity comparisons and absolute
bound used in the determinant extraction.

For clarity, consider the degree conditions in that extraction.
At an intervening position \(y\), let \(d_{\rm low}(y)\) be the
outer digit degree and \(d_A(y),d_G(y)\) the degrees of these two
factors. The class equation gives
\[
  d_{\rm low}(y)+p^{z_j-y}\bigl(d_A(y)+d_G(y)\bigr)\le\deg(y).
\]
Thus the separate order contributions can be paid out of their sum. At an incoming position
the only letter is \(W_{j-1,m}\), of degree \(b\), and the
complementary degree is zero.

Let \(R_{\rm low}\) be the original low radix and \(R_{\rm seg}\) the incoming segment radix at this letter. Substitution first multiplies the exponent by \(P_{j-1,m}\); passing to true degree then gives
\[
 b\ \longmapsto\ P_{j-1,m}b
      \ \longmapsto\ R_{\rm low}P_{j-1,m}b=R_{\rm seg}b .
\]
This is the contribution from the incoming letter alone. Other terms at the shared letter have nonnegative degree, so if \(d_{\rm coll}\) is its total collapsed plane degree and \(d_{\rm true}\) its true degree, then
\[
 R_{\rm seg}b\le R_{\rm low}d_{\rm coll}\le d_{\rm true}.
\] This argument applies
to every model monomial used, each of which has the prescribed
collapsed degrees. After collecting terms and projecting to a
segment multidegree the same upper bound holds. In particular
the upper allowance
\[
 1.01\deg(y)/p^{z_j-y}
\]
holds in both
places. The extraction proof in Section~\ref{sec:plane-lifts}
only uses this upper inequality when estimating the earlier rows.
It therefore applies also to components which have lost much of
their degree in the incoming frame. This is the one-sided form
of the realization lemma.

\proofstep{Recovering the original low coordinates.}
Multiply each product of segment factors by the original low
label after applying the segment power substitutions. The sum
is \(\operatorname{Asm}_\rho(\phi Z_*,N_0)\): the extra power
on an incoming row is exactly \(P_{j-1,m}\) and the label has
zero support off the tops. Within each fixed component the
result has the required full degrees.

In each true letter the substitutions multiply a segment degree by its radix. Summing over the contributing parts \(b\) therefore gives 
\[
 d_{\rm collapsed}=\sum_b P_b d_b,
\]
 where \(P_b,d_b\) are the radix and degree of the intervening, incoming or outgoing part. Restoring the label then gives the original degree
by the class equation. Its normal degree is also constant: at
an even top add its normal low digits to the low radix times the
complement quotient degree; at an odd top use its plane digits
and plane quotient degree. Sum these numbers over the tops.
If \(S\) is the component of order \(o\), the one-sided realization gives an integral \(P\) with
\[
 P|_{\rm chosen}\equiv w^o S\pmod{w^{o+c_P}},\qquad P|_{\rm competitor}\equiv0\pmod{w^{o+c_P}}
\]
at every competing plane. This proves the symbol membership and its realization on the stated one-sided range.

\proofstep{Keeping one order for all tested generators.}
Finally fix an order before choosing \(H\). The counted generators
have normal degree bounded by
\[
 X=N(\mathcal L a/(pk)+10a p^{-L_1}+kQ)
\]
by
Section~\ref{par:plane-rank}. The common factor has normal degree zero. For a counted generator \(A\) and a product of \(j\le Nc_0\) row operators \(D_1,\ldots,D_j\), the per-operator increment bound gives
\[
 \deg_{\rm normal}(D_1\cdots D_j A)
     \le \lfloor X\rfloor+j\le t_F,\qquad t_F=\lfloor X\rfloor+Nc_0,
\]
whenever the result is nonzero. Taking the span preserves this upper bound. The Fock form pairs only equal exponent tuples.

If \(A=\sum_{o\le t_F}A_o\) is a tested input and
\(Z=\sum_o Z_o\) is the assembled symbol, then the diagonal
Fock pairing gives
\[
                  [A,Z]=\sum_{o\le t_F}[A_o,Z_o] .
\]
The discarded terms have disjoint exponent support from \(A\).
After this deletion the ultrametric inequality selects a
summand with pairing valuation no greater than that of the sum. We have proved the flat symbol estimate
at the same \(D_X,t_F,c_P\) for every primitive polynomial in
the exact model span. \label{cor:plane-return-1}
\end{proof}


Since the evaluation on the sum has valuation at most \(D_X\),
one of its realizable summands of order \(o\le t_F\) has
valuation at most \(D_X\).

\paragraph{Returning to the restored weight.}

The flat symbol has now been lifted with uniform depth. We must compare its detecting covector with one for the original highest weight.

Write \(\lambda\) for the full weight, \(\mu\) for its
flat part and \(A_\eta=A_r(\eta)\) for the normalisation
exponent of a weight \(\eta\). In the notation of the
edge pairing, \(R^\eta_e=w^{-A_\eta}H^\eta_e\),
\(R^{\prime\eta}_e=w^{-A_\eta}(H^\eta_e)^*\), and the
normalised evaluation is \(b^\eta_e(u,T)=w^{2A_\eta}T(u)\).
\begin{established}[Transport maps and the uniform primitive-vector bound]\label{hyp:plane-full-weight-transport}
The operator construction needed for transport has the
following properties. Its maps \(\rho_\beta:V_\lambda\to
V_\mu\) send \(R^\lambda_e\) to \(R^\mu_e\). At every
other edge into the same target they lose at most
\[
 s=rNc_0+2rn'\mathrm{cpp}
\]
powers of \(w\) on the corresponding lattices. On the
target vertex lattices they map \(H^\lambda_+\) into
\(w^{2(A_\lambda-A_\mu)-s}H^\mu_+\). They send products
of a flat generator and a remaining-weight vector into
the allowed span of operator words on that generator.
Finally, for every primitive \(u\in R^\lambda_e\), some
\(\rho_\beta u\) lies outside \(w^{2\kappa+1}R^\mu_e\),
where \(\kappa\) is bounded uniformly before \(u\) is chosen.
\end{established}
\begin{proof}
Take the maps in Proposition~\ref{prop:plane-bounded-weight-operators}. Its divided-power construction at each tensor position preserves the selected edge lattice. The sum of the losses from the other positions is \(rNc_0+2rn'\mathrm{cpp}\). The integral dual estimate there finds an image of a primitive vector outside \(w^{2\kappa+1}R_e^\mu\); its bound on \(\kappa\) depends only on the weight spreads. The construction acts by words on the remaining tensor factors and hence has the required action on the span.
\end{proof}

\begin{proposition}[A flat covector transported to the full weight]\label{prop:plane-restore-covector}
Assume the lattice maps and primitive-vector bound of
Lemma~\ref{hyp:plane-full-weight-transport}. Suppose every
primitive flat vector in the prescribed operator span has a
flat covector of depth \(D_F\), endpoint order \(t_F\), and
competitor conductor \(c_P=c_\star+s\). Then every primitive
\(u\in R^\lambda_e\) in the span of the restored model family
has a full-weight covector whose depth \(D\), endpoint order \(o\) and competitor conductor \(c\) satisfy
\[
 D\le\kappa+D_F,\qquad o\le t_F+s,\qquad c=c_\star.
\]

\end{proposition}

\begin{proof}
\proofstep{Testing a primitive combination in the restored lattice.}
For a chosen primitive combination \(u\), select such a
\(\beta\) and write \(\rho_\beta u=w^a u_1\) with \(u_1\)
primitive and \(0\le a\le2\kappa\). The operator-word
span is over the fraction field, so it contains \(u_1\).
Transport a flat covector \(T\) by
\(\Lambda T=w^{2(A_\mu-A_\lambda)}\rho_\beta^* T\).
For example \(H^\lambda_e\) maps by \(\rho_\beta\)
into \(w^{A_\lambda-A_\mu}H^\mu_e\). Pulling back the
flat lattice \(w^{-A_\mu}(H^\mu_e)^*\) and inserting
the scalar \(w^{2(A_\mu-A_\lambda)}\) gives precisely
the required inclusion in \(R^{\prime\lambda}_e\).

The calculation on the competing edge lattices loses \(s\)
powers of \(w\); the stated vertex inclusion also gives
loss \(s\) at the target vertex. The normalised pairings obey
\[
       b^\lambda_e(u,\Lambda T)=w^a b^\mu_e(u_1,T).
\]
The pairing gains valuation \(a/2\le\kappa\), since
\(w^2=p\). At the vertex and the competing edges the lattice
bounds lose at most \(s\) powers of \(w\), while the chosen
edge retains integrality. Thus the transported test has
depth at most \(\kappa+D_F\), order at most \(t_F+s\)
and conductor at least \(c_P-s=c_\star\).  The model-to-clean
congruence at these depths was checked in
Section~\ref{par:plane-middle-weights}.

\end{proof}

The final step at a single edge is to place the detecting
functional in the target lattice of
Proposition~\ref{prop:residue-plane-neighbor-test}.
The dual-plane substitution accounts for the vanishing at
the other incoming edges.
\paragraph{Dual planes at the target.}
\label{par:plane-dual-edge}
\hypertarget{source-plane-edge-covector}{}
For clarity write the distinguished edge in adapted coordinates as
\(L_-=b W^k\), \(L_+=b D_p^{-1}W^k\), where
\(D_t=\operatorname{diag}(tI_r,I_{k-r})\). Its dual edge goes
from \(L_+^\vee=b^{-t}D_p W^k\) to
\(L_-^\vee=b^{-t}W^k\).

At each Frobenius position the middle
frame on this dual edge is the inverse transpose of the original
middle frame: in the original frame use \(bD_w^{-1}\), and in
the dual use
\[
 b^{-t}D_p D_w^{-1}=b^{-t}D_w,
\]
with the
indicated Frobenius action on \(b\). Distinct incoming edges at
\(L_+\) give distinct outgoing edges at \(L_+^\vee\) by double
duality. There they have distinct residue planes, since a lattice
between a lattice and its \(p^{-1}\)-multiple is determined by
its image in that quotient. For the standard target,
\(bD_p^{-1}\in\mathrm{GL}_k(W)\); it gives the Witt frame
whose inverse transpose is \(b^{-t}D_p\). Thus a single choice
of distinguished frame in the one-sided realization handles all
incoming competitors.

\begin{lemma}[A one-sided plane symbol gives an integral dual-edge test]\label{lem:plane-dual-edge-test}
For the edge from \(bW^k\) to \(bD_p^{-1}W^k\), use
\(bD_w^{-1}\) and \(b^{-t}D_w\) as the respective middle frames.
Let a detecting
symbol of order \(o\) be realized to precision \(o+c_P\)
at the distinguished dual edge and at least that order
at its competitors. If \(D_X<c_P/2\), the resulting flat covector on the original edge has order \(o'\), conductor \(c'\) and depth \(D'\) satisfying
\[
 o'\le o,\qquad c'\ge c_P,\qquad D'\le D_X.
\]

\end{lemma}

\begin{proof}
\proofstep{The dual pairing on the target plane.}
In flat polynomial coordinates use the bilinear form
\(p^{\kappa_{\rm fl}}[P,Q]/\prod_i d_i!\).
It is integral on coefficient lattices in inverse transpose
frames. Indeed the factorial deficiency gives this in the
standard coordinates; the substitution identity
\[
 [P(MX),Q(X)]=[P(X),Q(M^tX)]
\]
follows by writing the
pairing as a constant-coefficient differentiation at zero.
The identity applies to each letter, also when it carries the
inverse transpose representation. The scalar adjustments to the two edge substitutions give 
\[
 w^{A_\mu}w^{A_\mu}=w^{2A_\mu}=p^{A_\mu},
\]
 precisely
the normalized edge pairing for the flat weight \(\mu\), by
Section~\ref{par:plane-middle-weights}.

Take the integral polynomial realizing a detecting symbol
\(\sigma\) of order \(o\), and view it as a flat section in the
dual vertex frame. Divide it by \(w^o\), and pair by this form to
obtain \(T\). At the target its order is at most \(o\), by
integrality before division. In normalized coordinates on the
dual edge its polynomial after division is \(\sigma+w^{c_P}E\)
with integral \(E\); on any other dual edge at the same source
it belongs to \(w^{c_P}O[\mathbf X]\). The inverse transpose
calculation proves \(T\in R_e^{\prime\mu}\) and
\[
 T\in w^{c_P}R_{e'}^{\prime\mu}
\]
on every competing incoming
edge. For any normalized integral input, the model error pairs into
\(w^{c_P}O\), because the factorial rescaling is integral on
the paired edge lattices. Its valuation is at least
\[
 c_P/2,
\]

strictly larger than the detected valuation bounded by \(D_X\).
The valuation of the actual pairing therefore equals that
of the model detecting value.

\end{proof}

\paragraph{Operators for the bounded remaining weights.}
\label{par:plane-rho}
First retain only the top positions in the full weight \(\lambda\).
Write \(\lambda=\mu+\eta\) there, with \(\mu\) the flat weight
and \(\eta\) the remaining weight. Its spread at each top is at
most \(c_0\). We need two facts about the Schur lattices
of Section~\ref{sec:edge-contraction}.

First, for each of the weights \(\xi=\lambda,\mu\) there is a
contravariant identification \(\Psi_\xi:V_\xi^*\longrightarrow
V_\xi\) satisfying, simultaneously in every vertex frame \(g\),
\[
 p^{h_\xi} H_\xi(g^{-t})\ \subseteq \Psi_\xi\bigl(H_\xi(g)^*\bigr)\ \subseteq\ H_\xi(g^{-t}).
\]
In this display \(g\) denotes the tuple of the component frames
and \(H^*\) the integral dual. The same inclusions hold with
\(R_e^{\prime\xi},R_{e^\vee}^\xi\) in the middle and outside,
since the two normalizations multiply the unscaled lattices by
the same \(w^{-A_\xi}\).
Here \(h_\xi\) is the sum over positions of
\(\binom k2\lfloor\log_p(l+k-1)\rfloor\) at positive spreads
\(l\), with zero at spread zero.

\begin{proposition}[Uniform operators for the remaining weights]\label{prop:plane-bounded-weight-operators}
With the bounded top and off-top spreads, the maps
\(\rho_\beta\) of Proposition~\ref{prop:plane-restore-covector}
satisfy its three lattice conditions with
\(s=rNc_0+2rn'\mathrm{cpp}\). Every primitive vector
has an image outside \(w^{2\kappa+1}R_e^\mu\) for some
such map. On the restored model family in the prescribed range of parameters, the target depth \(D\), order \(o\) and conductor \(c\) satisfy
\[
 D\le D_X+\kappa\le54Nk^2\log_p a,\qquad o\le t_F+s,\qquad c=c_\star.
\]

\end{proposition}

\begin{proof}
\proofstep{Normalizing one position in the tensor.}
To verify the normalization at one position, remove the determinant
twist. Take the matrix-polynomial Fock form divided by the
squared Fock norm of the highest polynomial \(v\), the product
of the initial principal minors prescribed by the shape.
This norm is a positive integer. Pairing with \(v\) gives evaluation
at the identity: both vanish outside its weight by row and column
degrees, and both take value one on \(v\). The highest weight
space has dimension one. For a divided-power word applied to
\(v\), transpose the word in the pairing. It follows that its
pairing with any integral polynomial in the Schur lattice is
integral.

Conversely, integrality against the Schur lattice forces
\[
 Q\in O[\mathbf X]
\]
in the matrix coordinates.
Indeed substitute a generic matrix \(M\) in the highest polynomial
by the row action. The coefficients of the resulting polynomial
in the entries of \(M\) are integral elements of the Schur
lattice, by the minor expansion. Pair it with \(Q\) and transpose
the row action. The result is \(Q\) evaluated on the generic
transposed matrix, by the identity evaluation just proved.
Hence every coefficient of \(Q\) is integral, and saturation
of the standard monomial lattice gives the assertion.
The same argument proves nondegeneracy. The divided-power span of
\(v\) contains
\[
 p^{\binom k2\lfloor\log_p(l+k-1)\rfloor}
\]

times the Schur lattice by the trace argument of
Section~\ref{sec:edge-contraction}. Thus the integral dual for the
normalized form lies between the two asserted lattices. Restoring
the determinant twists on inverse transpose frames cancels their
two scalar factors. Tensoring the forms and integral dual bases
adds the exponents and gives \(\Psi_\xi\). Its inverse pairing
is a nonzero scalar times the matrix-polynomial Fock form.
Here \(l+k-1\le a^2\) for both weights at a sufficiently late
host: the flat degree is at most \(2a\), the residual spread
is at most \(c_0\), and the host size bound gives
\(2a+c_0+k\le a^2\). Consequently
\[
 \kappa=h_\lambda+h_\mu\le4L.
\]

Second, for residual sections \(s,\beta\) consider the flat
endomorphism \(U_{s,\beta}\) defined by
\[
  [\theta(U_{s,\beta}F),\theta H]_{\rm mat}
       =[\theta F\cdot s,\theta H\cdot\beta]_{\rm mat}.
\]
The matrix-polynomial form restricted to the image of the flat
model is nondegenerate, by the preceding argument. The flat
insertion \(\theta\) uses the first column or the cofactors.
The maps \(U_{s,\beta}\) are linear in either residual input.
At a single position their span is stable under commutators
with the row operators \(E_{ij}=X_i\partial_j\).
In fact transpose adjunction and the product rule move the
commutator to the two terms applying row operators to \(s\) and
to \(\beta\), with opposite signs and transposed indices.

For cofactors the induced row action on the flat coordinates is
the negative transpose action plus the scalar trace term, by
Jacobi's identity. If the induced operator is \(A+\lambda I\), its commutator on
an endomorphism \(U\) is
\[
 (A+\lambda I)U-U(A+\lambda I)=AU-UA .
\]
Thus the determinant twist introduces no extra row weight.
The residual space is generated by products of flag minors with
at most \(c_0\) columns. Each row occurs at most once in a
column. Thus the commutator weight of a generator
\(U_{s,\beta}\) has each coordinate between \(-c_0\) and \(c_0\):
it is the difference, in one of the two orders, of the two
residual row weights. The determinant twists cancel in this difference.

\proofstep{Operators from the diagonal action.}
For degree \(d\) the endomorphism representation of
\(\operatorname{Sym}^d\) in characteristic zero has one summand for each of the highest weights
\[
 \lambda^{(i)}=(i,0,\ldots,0,-i),\qquad 0\le i\le d.
\]
Indeed \(\operatorname{Sym}^d\otimes
(\operatorname{Sym}^d)^*\) is the bidegree \((d,d)\)
space decomposed by the powers of the contraction in
Section~\ref{sec:edge-transfer}; its harmonic quotients are
the irreducible incidence spaces described there. For summand
\(i\) the row word \(E_{1k}^i\) has that highest weight and
is killed under commutator by every raising operator. Its action on a single monomial gives
\[
 E_{1k}^i X_k^d
 =\frac{d!}{(d-i)!}X_1^i X_k^{\,d-i}\ne0
 \qquad(0\le i\le d),
\]
including the identity at \(i=0\). Words of length at
most \(i\) form a commutator-stable span by the product rule and
the matrix-unit commutator. They therefore contain the entire
summand.

Complete reducibility now implies that the span of the
\(U_{s,\beta}\) lies in words of length at most \(c_0\). Indeed
any summand present has its highest weight present, and the
weight bound just proved gives \(i\le c_0\).
At several positions expand the residual sections in pure
tensors and use factorization of the Fock form. Each term gives
at most \(c_0\) row operators at each top. Changes of flat frame
preserve this word space by conjugation. The word bound also holds when the flat degree is smaller than the allowed word length.

For \(\beta\in R_{e^\vee}^{\eta}\), write \(m_\beta\) for
multiplication by that residual section. Define
\[
   \Lambda_\beta T
       =p^{h_\lambda}\Psi_\lambda^{-1}
                      m_\beta\Psi_\mu T,\qquad
   T(\rho_\beta u)=p^{A_\eta}(\Lambda_\beta T)(u).
\]
The second equality defines a linear map by finite-dimensional
duality. The edge normalizations give
\[
 b_e^\mu(\rho_\beta u,T)=b_e^\lambda(u,\Lambda_\beta T).
\]
For \(T\in R_e^{\prime\mu}\), multiplication on the dual edge and
the inclusions for \(\Psi_\xi\) give
\[
 \Psi_\mu T\in R_{e^\vee}^\mu,\qquad
 m_\beta\Psi_\mu T\in R_{e^\vee}^\lambda,\qquad
 p^{h_\lambda}R_{e^\vee}^\lambda\subseteq
                              \Psi_\lambda(R_e^{\prime\lambda}).
\]
Thus \(\Lambda_\beta T\in R_e^{\prime\lambda}\).
In unscaled dual bases the two lattice factors \(w^{-A}\) cancel
the pairing factor \(w^{2A}\). The preceding adjoint identity
and this perfectness imply
\[
 \rho_\beta R_e^\lambda\subseteq R_e^\mu.
\]

On an input
\(u=\theta F\cdot s\), the definition of the two forms shows
that its flat output is a scalar multiple of
\(U_{s,\beta}F\) in polynomial coordinates. This proves
membership in the exact operator span.

\proofstep{Reservations for the remaining weights.}
Put \(s_T=rNc_0\). Multiplying any \(\beta\) from the chosen
dual edge by \(w^{s_T}\) puts it in the dual-edge source vertex
lattice. To see this, at each position the edge normalization
subtracts the top \(r\)-row weight sum. The largest possible
difference between that sum and a sum on any
\[
 r
\]
rows is at
most \(r\) times the weight spread, by the tableau weight bound.
This proves the claim in an adapted diagonal frame, including
determinant twists, and hence in every frame. Every outgoing normalized dual edge contains the vertex lattice.
Hence the cost outside the distinguished edge is at most \(s_T\).
For the target vertex, test against \(T\in(H_+^\mu)^*\). The vertex
inclusions and \(p=w^2\) give
\begin{equation}\label{eq:edge-target-adjoint}
 \begin{aligned}
 w^{s_T}\Lambda_\beta T &\in(H_+^\lambda)^*,\\
 T(\rho_\beta u)&=w^{2A_\eta-s_T}
          (w^{s_T}\Lambda_\beta T)(u)\in w^{2A_\eta-s_T}O
                      \qquad(u\in H_+^\lambda).
 \end{aligned}
\end{equation}
Testing every integral dual functional in \eqref{eq:edge-target-adjoint} yields
\(\rho_\beta H_+^\lambda\subseteq w^{2A_\eta-s_T}H_+^\mu\).

There is also a lower bound for this family of maps. The full
lattice on the dual edge is the integral span of products
\(y\beta\) from the flat and residual edge lattices. In a standard
tableau product, the column heights of the sum shape are the
union, with multiplicities, of those of the two shapes. Split
the columns accordingly and take a tensor of these factorizations
over the positions. The frame action, determinant twist and
edge powers all multiply. Thus for every full dual integral
functional \(G\), \(\Psi_\lambda G\) is an integral sum of these
products.

If
\[
 \Psi_\lambda G=\sum_i y_i\beta_i,
\]
take
\(T_i=\Psi_\mu^{-1}(p^{h_\mu}y_i)\), which is flat
dual integral. The multiplier in \(\Lambda_\beta\) then gives
\[
 p^\kappa G=\sum_i\Lambda_{\beta_i} T_i
\]
with
\(\kappa=h_\lambda+h_\mu\). For the normalized pairings,
\[
 p^\kappa b_e^\lambda(u,G)
       =\sum_i b_e^\mu(\rho_{\beta_i}u,T_i).
\]
Suppose \(\rho_\beta u\in w^{2\kappa+1}R_e^\mu\) for every
integral choice of \(\beta\). Testing the displayed sum against the
integral flat duals then gives, for every \(G\in R_e^{\prime\lambda}\),
\[
 p^\kappa b_e^\lambda(u,G)\in w^{2\kappa+1}O
 \quad\Longrightarrow\quad b_e^\lambda(u,G)\in wO
 \quad(p^\kappa=w^{2\kappa}).
\]
Perfectness gives \(u\in w R_e^\lambda\), a contradiction for a
primitive \(u\). Thus at least one residual test detects it.

\proofstep{Restoring the non-top tensor factor.}
Restore the non-top tensor factor, of weight \(\xi\), by taking
all contractions on it with \(p^{A_\xi}G\),
\(G\in R_e^{\prime\xi}\), before applying the top maps.
In a pair of normalized dual bases these contractions are
ordinary coordinate extractions. Some extraction of a primitive
tensor remains primitive. Thus this step adds no depth. The
spread off the tops is at most \(2\mathrm{cpp}\), so
\[
 w^{2rn'\mathrm{cpp}}G
\]
belongs to the target dual vertex
lattice and hence to all competing incoming dual edge lattices.

For this inclusion, the passage from an intermediate to a target
frame expands the first \(r\) directions once more by \(w^{-1}\);
a coordinate functional then loses at most the difference of two
\(r\)-row weight sums. On the target itself the contraction
carries the additional scalar \(p^{A_\xi}\). Adding the competing-edge loss from these positions to the
top contribution gives the budget
\[
 s=s_T+2rn'\mathrm{cpp},\qquad
 2A_\eta+2A_\xi-s=2(A_\lambda-A_\mu)-s .
\]
Here the right-hand exponent is the resulting target-vertex exponent.
The flat
degree is zero at every omitted position, so this contraction
retains the top word allowance.

Apply the adjoint calculation to the primitive restored model
combination. The three resulting bounds are
\begin{equation}\label{eq:restored-covector-budget}
 \begin{aligned}
 \text{endpoint order}&\le t_F+s\le\lceil t_{\rm II}\rceil,\\
 \text{depth}&\le D_X+\kappa\le54Nk^2\log_p a\le D_{\rm II},\\
 \text{competing conductor}&\ge c_P-s=c_\star,
                    \qquad c_\star=8(D_{\rm II}+1).
 \end{aligned}
\end{equation}
The last depth comparison and the endpoint bound in
\eqref{eq:restored-covector-budget} are the host choices of Section~\ref{sec:host-edge-tests},
with its absolute coefficient taken sufficiently large.
On the source side cleaning
the normal-degree components of the original flat generators
costs at most \(\lfloor X\rfloor\) at the source, by the
one-sided lift and division by the order of each component.
Multiplication by a remaining-weight section from the chosen
edge costs at most \(s\) at that vertex, by the same spread
estimate. Thus the source order also fits. With \(c'=c_P\)
one has \[
 c'\ge2(D_X+\kappa+1),\qquad c'\ge c_\star+s.
\]
Section~\ref{par:plane-middle-weights} transfers the pairing
to the clean combination and gives its competing source
conductors. This completes the single-edge estimate with the
order and depth fixed on the host and the stated quantifier over
primitive combinations. \label{cor:plane-edge-covector-1}
\end{proof}



\paragraph{From one plane to a full target.}
\begin{proof}[Conclusion of the proof of Proposition~\ref{prop:residue-plane-neighbor-test}]

Write \(d=\dim V\), and put \(a_0=a-c_G\).
The preceding construction on every edge into the standard
vertex gives \(m_e\) clean source vectors, with
\[
 m_e\ \ge\
 (1-CNr^2 k/Y_2)(1-k c_0/a_0)^N
                d\, q^{-r(k-r)} .
\]
For each combination within this block having a unit
coefficient, the covector of Section~\ref{par:plane-rho} evaluates
with valuation at most \(D_{\rm II}\) and satisfies the conductor
condition at every competing incoming edge. Both endpoint orders
are bounded by \(\lceil t_{\rm II}\rceil\).

Take incoming edges for which the codimension-\(r\)
subspace of \(L_+/pL_+\) is a graph over the last \(k-r\)
coordinates. All \(q^{r(k-r)}\) graphs give distinct lattices.
An arbitrary target is reached by a basis change over \(K_0\)
and its Frobenius twists. For a primitive combination of the
vectors from all these edges, choose a block containing a unit
coefficient and its detecting covector. The other blocks contribute
terms of valuation at least
\[
 c/2>D_{\rm II},
\]
by integrality
of the edge pairings and \(c=8(D_{\rm II}+1)\). They preserve
the valuation of the detected value.

\proofstep{Combining the independent directions at the target.}
If
\[
 M=\sum_e m_e,
\]
the evaluation image of the target lattice
in \(O^M\) contains \(p^{D_{\rm II}}O^M\). Indeed in Smith
coordinates any missing direction or any elementary divisor of
larger valuation gives a primitive linear combination taking
only values of larger valuation, contrary to what was just
proved. Choosing a preimage of each scaled basis vector gives
the integral right inverse at depth \(D_{\rm II}\). Finally
\[
 (1-kc_0/a_0)^N\ge 1-Nkc_0/a_0
\]
and \(c_0/a_0=O(1/X')\).
The resulting deficit is at most \(\min(1,C'Nr^2 k/Y_2)\)
after increasing the constant, including the range where this
bound is one. At \(Y_2=x^{26}\) this deficit is \(O(Nx^{-12})\)
with the host bounds of Section~\ref{sec:host-edge-tests}.
The finite extension argument there preserves this estimate
on the selected good cycles. This proves the smaller-prime input
at the same host as the larger-prime estimate.
 \label{cor:plane-lifts-1}
\end{proof}
